% !TeX program = lualatex
% ================================================================
% Polish Parallel Text Version of FactorsMultiplesPrimesQuickQuestions
%
% Generated by the server using the following settings:
%
%   language            Polish (pol)
%   text direction      left to right
%   font                Noto Serif
%   fallback font       Noto Serif
%   built from          latex/quickQuestions/factorsMultiplesPrimesQuickQuestions.tex
%   vocabulary key      latex/quickQuestions/factorsMultiplesPrimesQuickQuestions/L2/pol/factorsMultiplesPrimesQuickQuestions_polishKey.tex
%   tier 3 vocabulary   green   RGB 0 128 0
%   English text        black   RGB 0 0 0
%   Polish text         purple  RGB 112 48 160
%   layout macros       version 3
%
% Please don't edit any information in this comment block - the server
% compares it against the current settings and rebuilds this file when
% they differ, so changes here would be overwritten. However, feel free
% to make updates to the LaTeX below if you want to edit it.
% ================================================================

\documentclass[aspectratio=169]{beamer}
% ================================================================
% Copied from the English sheet this was translated from, so that
% anything it sets up for itself still works here
% ================================================================
\usepackage{amsmath}
\usepackage{amssymb}
\usepackage{tikz}
\usetikzlibrary{shapes.geometric}

% ================================================================
% Quick Questions deck: factors, multiples and primes.
%
% Each topic opens with five true/false statements and then runs ten
% quick questions that get gradually harder. Every question gets two
% slides: the question on its own for the mini whiteboards, then the
% same question again with the solution underneath in red.
%
% Two topics are built differently. "Is it prime?" runs every whole
% number from 1 to 100 on its own slide, then lists that number's
% factors and says why it is or is not prime. "HCF and LCM from lists"
% builds each answer over five slides, adding the factor lists, the
% multiples and then the rings round the HCF and the LCM one slide at
% a time.
% ================================================================

\setbeamertemplate{navigation symbols}{}

% A link in the bottom left of every slide back to the contents, so a
% deck this long can be jumped around in during a lesson.
\setbeamertemplate{footline}{%
  \hbox to\paperwidth{%
    \hspace{1em}%
    \hyperlink{contents}{%
      \color{gray}\scriptsize$\leftarrow$~Back to contents}%
    \hfil}%
  \vskip4pt%
}

% Topic divider.
\newcommand{\qtopic}[1]{%
  \section{#1}%
  \begin{frame}
    \vfill
    \begin{center}
      {\usebeamercolor[fg]{structure}\Huge\bfseries #1}
    \end{center}
    \vfill
  \end{frame}%
}

% \qq[<question size>][<solution size>]{<question>}{<solution>}
% The two optional sizes drop the type for the wordier questions and
% the longer solutions, so nothing runs off the slide.
\NewDocumentCommand{\qq}{O{\Huge}O{\LARGE}+m+m}{%
  \begin{frame}
    \vfill
    \begin{center}
      \begin{minipage}{0.92\textwidth}
        \centering #1 #3
      \end{minipage}
    \end{center}
    \vfill
  \end{frame}%
  \begin{frame}
    \vfill
    \begin{center}
      \begin{minipage}{0.92\textwidth}
        \centering #1 #3
        \par\vspace{0.9em}
        {\color{red}#2 #4}
      \end{minipage}
    \end{center}
    \vfill
  \end{frame}%
}

% \tf[<statement size>][<answer size>]{<statement>}{<answer>}
% As \qq, but with a standing "True or false?" heading above the
% statement. The statements hunt for the usual misconceptions, so the
% answer always says why, not just which.
\NewDocumentCommand{\tf}{O{\LARGE}O{\Large}+m+m}{%
  \begin{frame}
    \vfill
    \begin{center}
      \begin{minipage}{0.92\textwidth}
        \centering
        {\usebeamercolor[fg]{structure}\Large\bfseries True or false?}
        \par\vspace{1em}
        #1 #3
      \end{minipage}
    \end{center}
    \vfill
  \end{frame}%
  \begin{frame}
    \vfill
    \begin{center}
      \begin{minipage}{0.92\textwidth}
        \centering
        {\usebeamercolor[fg]{structure}\Large\bfseries True or false?}
        \par\vspace{1em}
        #1 #3
        \par\vspace{0.9em}
        {\color{red}#2 #4}
      \end{minipage}
    \end{center}
    \vfill
  \end{frame}%
}

% \prq{<number>}{<factors>}{<why it is or is not prime>}
% The "Is it prime?" drill: the number alone in a large centred font,
% then the same number with all of its factors and the reason in red.
% The number is scaled rather than set with \fontsize, because the
% sans font stops at 51.59pt and larger sizes are silently substituted.
\NewDocumentCommand{\prq}{m m +m}{%
  \begin{frame}
    \vfill
    \begin{center}
      \scalebox{3.6}{\Huge\bfseries #1}
    \end{center}
    \vfill
  \end{frame}%
  \begin{frame}
    \vfill
    \begin{center}
      \scalebox{1.9}{\Huge\bfseries #1}
      \par\vspace{0.8em}
      \begin{minipage}{0.9\textwidth}
        \centering\color{red}
        {\large Factors of $#1$: #2}
        \par\vspace{0.7em}
        {\large #3}
      \end{minipage}
    \end{center}
    \vfill
  \end{frame}%
}

% Rings drawn round a number in a list. The ring is set in a box the
% width of the number itself, so the list is spaced as though it were
% ordinary text and nothing moves on the slide where the ring appears;
% the ring itself spills over the gaps either side, as a ring drawn by
% hand would.
\tikzset{
  ringed/.style={ellipse, inner sep=1pt, draw=red, line width=0.9pt},
  unringed/.style={ellipse, inner sep=1pt, draw=none},
}
\newlength{\ringwd}
\newcommand{\ringnode}[2]{%
  \settowidth{\ringwd}{$#2$}%
  \addtolength{\ringwd}{3pt}% a little clearance for the ring to sit in
  \makebox[\ringwd][c]{%
    \tikz[baseline=(n.base)]{\node[#1] (n) {$#2$};}%
  }%
}
\newcommand{\hcfmark}[1]{\alt<4->{\ringnode{ringed}{#1}}{\ringnode{unringed}{#1}}}
\newcommand{\lcmmark}[1]{\alt<5->{\ringnode{ringed}{#1}}{\ringnode{unringed}{#1}}}

% \hcflcm{<a>}{<b>}{<factors of a>}{<factors of b>}
%        {<multiples of a>}{<multiples of b>}
% Five slides: the question, the two factor lists, the two lists of
% multiples, the ring round the HCF, then the ring round the LCM.
% Everything added after the question slide is in red.
% The frame is top aligned so that the question stays put as the lists
% are added underneath it.
\NewDocumentCommand{\hcflcm}{m m +m +m +m +m}{%
  \begin{frame}<1-5>[t]
    \vspace{0.8em}
    \begin{center}
      {\LARGE Find the HCF and LCM of $#1$ and $#2$.}
    \end{center}
    \vspace{1.6em}
    \begin{columns}[T]
      \begin{column}{0.40\textwidth}
        \only<2->{\color{red}\large
          Factors of $#1$:\\[0.35em]
          #3
          \par\vspace{1em}
          Factors of $#2$:\\[0.35em]
          #4}
      \end{column}
      \begin{column}{0.58\textwidth}
        \only<3->{\color{red}\large
          Multiples of $#1$:\\[0.35em]
          #5
          \par\vspace{1em}
          Multiples of $#2$:\\[0.35em]
          #6}
      \end{column}
    \end{columns}
    \vfill
  \end{frame}%
}

% Contents-slide bullets. Green then orange sets the difficulty band;
% blue marks the two topics that work through a method a slide at a
% time rather than asking for a quick answer.
\newcommand{\tocbullet}[1]{\textcolor{#1}{$\bullet$}}
\setbeamertemplate{section in toc}{%
  \ifcase\inserttocsectionnumber
    \or\tocbullet{green!55!black}% 1 Factors
    \or\tocbullet{green!55!black}% 2 Multiples
    \or\tocbullet{green!55!black}% 3 Primes
    \or\tocbullet{blue}%          4 Is it prime?
    \or\tocbullet{orange}%        5 Prime factor decomposition
    \or\tocbullet{blue}%          6 HCF and LCM from lists
    \or\tocbullet{orange}%        7 HCF and LCM by prime factor decomposition
    \or\tocbullet{orange}%        8 Square numbers from prime factors
    \or\tocbullet{orange}%        9 Problem solving with HCF and LCM
  \fi
  ~\inserttocsection%
}

\title{Factors, Multiples and Primes: Quick Questions}
\subtitle{Mini whiteboard questions, with solutions}
\author{}
\date{}

% Characters this language's font has not got are borrowed from another,
% rather than being silently dropped - see docs/translated-sheet-latex.md
\directlua{%
  if luaotfload and luaotfload.add_fallback then
    luaotfload.add_fallback("eallatinfallback",
      { "Noto Serif:mode=harf;" })
  else
    tex.error("this luaotfload is too old to register a fallback font")
  end
}
\usepackage[english]{babel}
\babelprovide[import]{polish}
\babelfont[polish]{rm}[Renderer=Harfbuzz, BoldFont={Noto Serif}, ItalicFont={Noto Serif}, BoldItalicFont={Noto Serif}, RawFeature={fallback=eallatinfallback}]{Noto Serif}
\babelfont[polish]{sf}[Renderer=Harfbuzz, BoldFont={Noto Serif}, ItalicFont={Noto Serif}, BoldItalicFont={Noto Serif}, RawFeature={fallback=eallatinfallback}]{Noto Serif}
\newcommand{\ealtext}[1]{\foreignlanguage{polish}{#1}}
\newcommand{\ealtextblock}[1]{{\raggedright\foreignlanguage{polish}{#1}\par}}
% ================================================================
% EAL layout helpers
% ================================================================
\usepackage{xcolor}

\definecolor{ealkeycolour}{RGB}{0,128,0}
\definecolor{ealenglish}{RGB}{0,0,0}
\definecolor{ealtwocolour}{RGB}{112,48,160}

\newsavebox{\ealboxen}
\newsavebox{\ealboxtr}
\newlength{\ealglosswd}
\newlength{\ealglossraise}
\setlength{\ealglossraise}{1.15em}

% a tier 3 word, in English and in translation alike
\newcommand{\ealkey}[1]{\textcolor{ealkeycolour}{#1}}
\newcommand{\ealkeytr}[1]{\textcolor{ealkeycolour}{\ealtext{#1}}}

% One translated word above one English word. Built from kernel
% primitives, never a tabular, which would break wherever the sheet
% loads array, booktabs or tabularx. The box is as wide as the wider of
% the two, so a long translation pushes its neighbours aside instead of
% overlapping them, and it claims its full height so a table row or a
% TikZ node makes room for it.
\newcommand{\ealgloss}[2]{%
  \sbox{\ealboxen}{\ealkey{#1}}%
  \sbox{\ealboxtr}{\small\textcolor{ealtwocolour}{\ealtext{#2}}}%
  \setlength{\ealglosswd}{\wd\ealboxen}%
  \ifdim\wd\ealboxtr>\ealglosswd\setlength{\ealglosswd}{\wd\ealboxtr}\fi
  \leavevmode
  \makebox[\ealglosswd][c]{%
    \makebox[0pt][c]{\usebox{\ealboxen}}%
    \makebox[0pt][c]{\raisebox{\ealglossraise}{\usebox{\ealboxtr}}}%
  }%
}

% opens up line spacing around a block of glosses, so the raised
% translations sit clear of the line above
\newenvironment{ealglossed}{\par\linespread{2.1}\selectfont\ignorespaces}{\par}

% a whole translated sentence above its English counterpart. block level,
% so it does not belong inside a tabular cell
\newcommand{\ealpara}[2]{%
  \par\addvspace{0.5\baselineskip}%
  {\color{ealtwocolour}\ealtextblock{#1}}%
  \penalty200
  {\color{ealenglish}#2\par}%
  \addvspace{0.5\baselineskip}%
}

\begin{document}

\frame{\titlepage}

\begin{frame}[label=contents]{Contents}
  \ealpara{Spis treści}{Contents}
  \small
  \tableofcontents
\end{frame}

%================================================================
\section{Factors}
\begin{frame}
  \vfill
  \begin{center}
    \ealpara{\ealkeytr{Dzielniki}}{\ealkey{Factors}}
    \par\vspace{1em}
    {\usebeamercolor[fg]{structure}\Huge\bfseries \ealkey{Factors}}
  \end{center}
  \vfill
\end{frame}
%================================================================

\tf{\ealpara{$7$ jest \ealkeytr{dzielnikiem} $21$}{$7$ is a \ealkey{factor} of $21$}}
  {\ealpara{\textbf{Prawda.} $21=7\times 3$, więc $21$ dzieli się dokładnie przez $7$.}{\textbf{True.} $21=7\times 3$, so $21$ divides exactly by $7$.}}

\tf{\ealpara{$1$ jest \ealkeytr{dzielnikiem} każdej liczby całkowitej}{$1$ is a \ealkey{factor} of every whole number}}
  {\ealpara{\textbf{Prawda.} Każdą liczbę $n$ można zapisać jako $1\times n$.}{\textbf{True.} Every number $n$ can be written as $1\times n$.}}

\tf{\ealpara{$6$ jest \ealkeytr{dzielnikiem} $16$}{$6$ is a \ealkey{factor} of $16$}}
  {\ealpara{\textbf{Fałsz.} \ealkeytr{Dzielniki} $16$ to $1,2,4,8,16$;\\[0.3em]
   $16\div 6$ nie jest liczbą całkowitą.}{\textbf{False.} The \ealkey{factors} of $16$ are $1,2,4,8,16$;\\[0.3em]
   $16\div 6$ is not a whole number.}}

\tf{\ealpara{Każda liczba ma parzystą liczbę \ealkeytr{dzielników}}{Every number has an even number of \ealkey{factors}}}
  {\ealpara{\textbf{Fałsz.} \ealkeytr{Dzielniki} zwykle łączą się w pary, ale $9$ ma tylko $1,3,9$ --- trzy \ealkeytr{dzielniki} --- ponieważ $3\times 3$ używa tego samego \ealkeytr{dzielnika} dwa razy.}{\textbf{False.} \ealkey{Factors} usually pair up, but $9$ has just $1,3,9$
   --- three \ealkey{factors} --- because $3\times 3$ uses the same \ealkey{factor} twice.}}

\tf{\ealpara{Jeśli $4$ jest \ealkeytr{dzielnikiem} liczby,\\[0.3em]to $8$ też musi być jej \ealkeytr{dzielnikiem}}{If $4$ is a \ealkey{factor} of a number,\\[0.3em]then $8$ must be a \ealkey{factor} of it too}}
  {\ealpara{\textbf{Fałsz.} $4$ jest \ealkeytr{dzielnikiem} $12$, ale $8$ nie jest.}{\textbf{False.} $4$ is a \ealkey{factor} of $12$, but $8$ is not.}}

\qq[\LARGE]{\ealpara{Wypisz wszystkie \ealkeytr{dzielniki} $12$}{List all the \ealkey{factors} of $12$}}
  {\ealpara{$1,\ 2,\ 3,\ 4,\ 6,\ 12$}{$1,\ 2,\ 3,\ 4,\ 6,\ 12$}}

\qq[\LARGE]{\ealpara{Wypisz wszystkie \ealkeytr{dzielniki} $20$}{List all the \ealkey{factors} of $20$}}
  {\ealpara{$1,\ 2,\ 4,\ 5,\ 10,\ 20$}{$1,\ 2,\ 4,\ 5,\ 10,\ 20$}}

\qq[\LARGE]{\ealpara{Wypisz wszystkie \ealkeytr{dzielniki} $36$}{List all the \ealkey{factors} of $36$}}
  {\ealpara{$1,\ 2,\ 3,\ 4,\ 6,\ 9,\ 12,\ 18,\ 36$}{$1,\ 2,\ 3,\ 4,\ 6,\ 9,\ 12,\ 18,\ 36$}}

\qq[\LARGE]{\ealpara{Ile \ealkeytr{dzielników} ma $16$?}{How many \ealkey{factors} does $16$ have?}}
  {\ealpara{$1,\ 2,\ 4,\ 8,\ 16$\\[0.3em]Pięć \ealkeytr{dzielników}.}{$1,\ 2,\ 4,\ 8,\ 16$\\[0.3em]Five \ealkey{factors}.}}

\qq[\LARGE]{\ealpara{Znajdź \ealkeytr{największy wspólny dzielnik} $12$ i $18$}{Find the \ealkey{highest common factor} of $12$ and $18$}}
  {\ealpara{\ealkeytr{Dzielniki} $12$: $1,\ 2,\ 3,\ 4,\ \mathbf{6},\ 12$\\[0.3em]
   \ealkeytr{Dzielniki} $18$: $1,\ 2,\ 3,\ \mathbf{6},\ 9,\ 18$\\[0.3em]
   $\mathrm{HCF}=6$}{\ealkey{Factors} of $12$: $1,\ 2,\ 3,\ 4,\ \mathbf{6},\ 12$\\[0.3em]
   \ealkey{Factors} of $18$: $1,\ 2,\ 3,\ \mathbf{6},\ 9,\ 18$\\[0.3em]
   $\mathrm{HCF}=6$}}

\qq[\LARGE]{\ealpara{Znajdź \ealkeytr{największy wspólny dzielnik} $24$ i $36$}{Find the \ealkey{highest common factor} of $24$ and $36$}}
  {\ealpara{\ealkeytr{Dzielniki} $24$: $1,\ 2,\ 3,\ 4,\ 6,\ 8,\ \mathbf{12},\ 24$\\[0.3em]
   \ealkeytr{Dzielniki} $36$: $1,\ 2,\ 3,\ 4,\ 6,\ 9,\ \mathbf{12},\ 18,\ 36$\\[0.3em]
   $\mathrm{HCF}=12$}{\ealkey{Factors} of $24$: $1,\ 2,\ 3,\ 4,\ 6,\ 8,\ \mathbf{12},\ 24$\\[0.3em]
   \ealkey{Factors} of $36$: $1,\ 2,\ 3,\ 4,\ 6,\ 9,\ \mathbf{12},\ 18,\ 36$\\[0.3em]
   $\mathrm{HCF}=12$}}

\qq[\LARGE]{\ealpara{Wypisz wszystkie \ealkeytr{pary czynników} $48$}{List all the \ealkey{factor pairs} of $48$}}
  {\ealpara{$1\times 48\qquad 2\times 24\qquad 3\times 16$\\[0.4em]
   $4\times 12\qquad 6\times 8$}{$1\times 48\qquad 2\times 24\qquad 3\times 16$\\[0.4em]
   $4\times 12\qquad 6\times 8$}}

\qq[\LARGE]{\ealpara{Która liczba całkowita poniżej $30$ ma najwięcej \ealkeytr{dzielników}?}{Which whole number below $30$ has the most \ealkey{factors}?}}
  {\ealpara{$24$, z ośmioma \ealkeytr{dzielnikami}:\\[0.3em]
   $1,\ 2,\ 3,\ 4,\ 6,\ 8,\ 12,\ 24$}{$24$, with eight \ealkey{factors}:\\[0.3em]
   $1,\ 2,\ 3,\ 4,\ 6,\ 8,\ 12,\ 24$}}

\qq[\LARGE]{\ealpara{$n$ ma dokładnie trzy \ealkeytr{dzielniki}.\\[0.4em]
Podaj możliwą wartość $n$.}{$n$ has exactly three \ealkey{factors}.\\[0.4em]
Give a possible value of $n$.}}
  {\ealpara{$4$ ($1,2,4$), $9$ ($1,3,9$) lub $25$ ($1,5,25$).\\[0.3em]
   Każda \ealkeytr{liczba pierwsza} podniesiona do kwadratu działa.}{$4$ ($1,2,4$), $9$ ($1,3,9$) or $25$ ($1,5,25$).\\[0.3em]
   Any \ealkey{prime} squared works.}}

\qq[\Large][\Large]{\ealpara{Wyjaśnij, dlaczego \ealkeytr{liczby kwadratowe} mają nieparzystą liczbę \ealkeytr{dzielników}}{Explain why \ealkey{square numbers} have an odd number of \ealkey{factors}}}
  {\ealpara{\ealkeytr{Dzielniki} występują w parach, po jednym z każdej strony \ealkeytr{pierwiastka}.\\[0.3em]
   W \ealkeytr{liczbie kwadratowej} \ealkeytr{pierwiastek} łączy się sam ze sobą,
   więc jeden \ealkeytr{dzielnik} zostaje bez pary.}{\ealkey{Factors} come in pairs, one either side of the \ealkey{square root}.\\[0.3em]
   In a \ealkey{square number} the \ealkey{square root} pairs with itself,
   so one \ealkey{factor} is left unpaired.}}

%================================================================
\section{Multiples}
\begin{frame}
  \vfill
  \begin{center}
    \ealpara{\ealkeytr{Wielokrotności}}{\ealkey{Multiples}}
    \par\vspace{1em}
    {\usebeamercolor[fg]{structure}\Huge\bfseries \ealkey{Multiples}}
  \end{center}
  \vfill
\end{frame}
%================================================================

\tf{\ealpara{$24$ jest \ealkeytr{wielokrotnością} $6$}{$24$ is a \ealkey{multiple} of $6$}}
  {\ealpara{\textbf{Prawda.} $6\times 4=24$.}{\textbf{True.} $6\times 4=24$.}}

\tf{\ealpara{Każda \ealkeytr{wielokrotność} $10$ jest także \ealkeytr{wielokrotnością} $5$}{Every \ealkey{multiple} of $10$ is also a \ealkey{multiple} of $5$}}
  {\ealpara{\textbf{Prawda.} $10=5\times 2$, więc każda \ealkeytr{wielokrotność} $10$
   już zawiera \ealkeytr{dzielnik} $5$.}{\textbf{True.} $10=5\times 2$, so every \ealkey{multiple} of $10$
   already contains a \ealkey{factor} of $5$.}}

\tf{\ealpara{$35$ jest \ealkeytr{wielokrotnością} $4$}{$35$ is a \ealkey{multiple} of $4$}}
  {\ealpara{\textbf{Fałsz.} $4\times 8=32$ i $4\times 9=36$, więc $35$ jest pominięte.}{\textbf{False.} $4\times 8=32$ and $4\times 9=36$, so $35$ is stepped over.}}

\tf{\ealpara{Każda \ealkeytr{wielokrotność} $3$ jest także \ealkeytr{wielokrotnością} $9$}{Every \ealkey{multiple} of $3$ is also a \ealkey{multiple} of $9$}}
  {\ealpara{\textbf{Fałsz.} $6$ jest \ealkeytr{wielokrotnością} $3$, ale nie $9$.\\[0.3em]
   W drugą stronę \emph{jest} to prawda.}{\textbf{False.} $6$ is a \ealkey{multiple} of $3$ but not of $9$.\\[0.3em]
   The other way round \emph{is} true.}}

\tf[\Large]{\ealpara{\ealkeytr{Najmniejsza wspólna wielokrotność} dwóch liczb jest zawsze znajdowana przez
pomnożenie ich przez siebie}{\ealkey{The lowest common multiple} of two numbers is always found by
multiplying them together}}
  {\ealpara{\textbf{Fałsz.} $4\times 6=24$, ale \ealkeytr{LCM} $4$ i $6$ to $12$.\\[0.3em]
   Mnożenie działa tylko wtedy, gdy liczby nie mają wspólnych \ealkeytr{dzielników}.}{\textbf{False.} $4\times 6=24$, but the \ealkey{LCM} of $4$ and $6$ is $12$.\\[0.3em]
   Multiplying only works when the numbers share no \ealkey{factors}.}}

\qq[\LARGE]{\ealpara{Zapisz pierwsze pięć \ealkeytr{wielokrotności} $4$}{Write down the first five \ealkey{multiples} of $4$}}
  {\ealpara{$4,\ 8,\ 12,\ 16,\ 20$}{$4,\ 8,\ 12,\ 16,\ 20$}}

\qq[\LARGE]{\ealpara{Zapisz pierwsze pięć \ealkeytr{wielokrotności} $7$}{Write down the first five \ealkey{multiples} of $7$}}
  {\ealpara{$7,\ 14,\ 21,\ 28,\ 35$}{$7,\ 14,\ 21,\ 28,\ 35$}}

\qq[\LARGE]{\ealpara{Zapisz \ealkeytr{wielokrotności} $6$ między $20$ a $50$}{Write down the \ealkey{multiples} of $6$ between $20$ and $50$}}
  {\ealpara{$24,\ 30,\ 36,\ 42,\ 48$}{$24,\ 30,\ 36,\ 42,\ 48$}}

\qq[\LARGE]{\ealpara{Czy $126$ jest \ealkeytr{wielokrotnością} $9$?}{Is $126$ a \ealkey{multiple} of $9$?}}
  {\ealpara{Tak. $1+2+6=9$, a $9\times 14=126$.}{Yes. $1+2+6=9$, and $9\times 14=126$.}}

\qq[\LARGE]{\ealpara{Znajdź \ealkeytr{najmniejszą wspólną wielokrotność} $4$ i $6$}{Find the \ealkey{lowest common multiple} of $4$ and $6$}}
  {\ealpara{$4,\ 8,\ \mathbf{12},\ 16,\ \ldots$\\[0.3em]
   $6,\ \mathbf{12},\ 18,\ \ldots$\\[0.3em]
   $\mathrm{LCM}=12$}{$4,\ 8,\ \mathbf{12},\ 16,\ \ldots$\\[0.3em]
   $6,\ \mathbf{12},\ 18,\ \ldots$\\[0.3em]
   $\mathrm{LCM}=12$}}

\qq[\LARGE]{\ealpara{Znajdź \ealkeytr{najmniejszą wspólną wielokrotność} $8$ i $12$}{Find the \ealkey{lowest common multiple} of $8$ and $12$}}
  {\ealpara{$8,\ 16,\ \mathbf{24},\ 32,\ \ldots$\\[0.3em]
   $12,\ \mathbf{24},\ 36,\ \ldots$\\[0.3em]
   $\mathrm{LCM}=24$}{$8,\ 16,\ \mathbf{24},\ 32,\ \ldots$\\[0.3em]
   $12,\ \mathbf{24},\ 36,\ \ldots$\\[0.3em]
   $\mathrm{LCM}=24$}}

\qq[\LARGE]{\ealpara{Znajdź \ealkeytr{najmniejszą wspólną wielokrotność} $5$, $6$ i $9$}{Find the \ealkey{lowest common multiple} of $5$, $6$ and $9$}}
  {\ealpara{\ealkeytr{LCM} $5$ i $6$ to $30$.\\[0.3em]
   \ealkeytr{LCM} $30$ i $9$ to $90$.\\[0.3em]
   $\mathrm{LCM}=90$}{\ealkey{LCM} of $5$ and $6$ is $30$.\\[0.3em]
   \ealkey{LCM} of $30$ and $9$ is $90$.\\[0.3em]
   $\mathrm{LCM}=90$}}

\qq[\LARGE]{\ealpara{Liczba jest \ealkeytr{wielokrotnością} zarówno $4$, jak i $10$.\\[0.4em]
Jaka jest najmniejsza możliwa liczba?}{A number is a \ealkey{multiple} of both $4$ and $10$.\\[0.4em]
What is the smallest it could be?}}
  {\ealpara{$20$, \ealkeytr{najmniejsza wspólna wielokrotność} $4$ i $10$.}{$20$, the \ealkey{lowest common multiple} of $4$ and $10$.}}

\qq[\Large][\Large]{\ealpara{Dwie latarnie morskie błyskają razem teraz.\\[0.4em]
Jedna błyska co $12$ sekund, druga co $18$ sekund.\\[0.4em]
Kiedy następnym razem błysną razem?}{Two lighthouses flash together now.\\[0.4em]
One flashes every $12$ seconds, the other every $18$ seconds.\\[0.4em]
When do they next flash together?}}
  {\ealpara{$\mathrm{LCM}(12,18)=36$\\[0.3em]
   Za $36$ sekund.}{$\mathrm{LCM}(12,18)=36$\\[0.3em]
   In $36$ seconds.}}

\qq[\Large][\Large]{\ealpara{Znajdź najmniejszą liczbę większą od $1$, która przy dzieleniu przez $4$, przez $5$ i przez $6$ daje resztę $1$}{Find the smallest number above $1$ that leaves a remainder
of $1$ when divided by $4$, by $5$ and by $6$}}
  {\ealpara{Liczba musi być o jeden większa od wspólnej \ealkeytr{wielokrotności} $4$, $5$ i $6$.\\[0.3em]
   $\mathrm{LCM}(4,5,6)=60$, więc liczba to $61$.}{The number must be one more than a common \ealkey{multiple} of $4$, $5$ and $6$.\\[0.3em]
   $\mathrm{LCM}(4,5,6)=60$, so the number is $61$.}}

%================================================================
\section{Primes}
\begin{frame}
  \vfill
  \begin{center}
    \ealpara{\ealkeytr{Liczby pierwsze}}{\ealkey{Primes}}
    \par\vspace{1em}
    {\usebeamercolor[fg]{structure}\Huge\bfseries \ealkey{Primes}}
  \end{center}
  \vfill
\end{frame}
%================================================================

\tf{\ealpara{$1$ jest \ealkeytr{liczbą pierwszą}}{$1$ is a \ealkey{prime} number}}
  {\ealpara{\textbf{Fałsz.} \ealkeytr{Liczba pierwsza} ma dokładnie dwa \ealkeytr{dzielniki}, ale $1$ ma tylko jeden.}{\textbf{False.} A \ealkey{prime} has exactly two \ealkey{factors}, but $1$ has only one.}}

\tf{\ealpara{$2$ jest jedyną parzystą \ealkeytr{liczbą pierwszą}}{$2$ is the only even \ealkey{prime} number}}
  {\ealpara{\textbf{Prawda.} Każda inna liczba parzysta ma $2$ jako dodatkowy \ealkeytr{dzielnik},
   więc ma więcej niż dwa \ealkeytr{dzielniki}.}{\textbf{True.} Every other even number has $2$ as an extra \ealkey{factor},
   so it has more than two \ealkey{factors}.}}

\tf{\ealpara{$51$ jest \ealkeytr{liczbą pierwszą}}{$51$ is a \ealkey{prime} number}}
  {\ealpara{\textbf{Fałsz.} $51=3\times 17$.\\[0.3em]
   Jej cyfry sumują się do $6$, więc dzieli się przez $3$.}{\textbf{False.} $51=3\times 17$.\\[0.3em]
   Its digits add to $6$, so it divides by $3$.}}

\tf{\ealpara{Istnieje nieskończenie wiele \ealkeytr{liczb pierwszych}}{There are infinitely many \ealkey{prime} numbers}}
  {\ealpara{\textbf{Prawda.} \ealkeytr{Liczby pierwsze} nigdy się nie kończą,
   choć stają się rzadsze, gdy liczby rosną.}{\textbf{True.} The \ealkey{primes} never run out,
   although they do thin out as the numbers get larger.}}

\tf{\ealpara{Jeśli liczba kończy się na $3$, to musi być \ealkeytr{liczbą pierwszą}}{If a number ends in $3$, it must be \ealkey{prime}}}
  {\ealpara{\textbf{Fałsz.} $33=3\times 11$ i $63=7\times 9$.}{\textbf{False.} $33=3\times 11$ and $63=7\times 9$.}}

\qq[\LARGE]{\ealpara{Zapisz wszystkie \ealkeytr{liczby pierwsze} między $1$ a $10$}{Write down all the \ealkey{prime} numbers between $1$ and $10$}}
  {\ealpara{$2,\ 3,\ 5,\ 7$}{$2,\ 3,\ 5,\ 7$}}

\qq[\LARGE]{\ealpara{Zapisz wszystkie \ealkeytr{liczby pierwsze} między $10$ a $20$}{Write down all the \ealkey{prime} numbers between $10$ and $20$}}
  {\ealpara{$11,\ 13,\ 17,\ 19$}{$11,\ 13,\ 17,\ 19$}}

\qq[\LARGE]{\ealpara{Czy $27$ jest \ealkeytr{liczbą pierwszą}?\\[0.4em]Wyjaśnij swoją odpowiedź.}{Is $27$ a \ealkey{prime} number?\\[0.4em]Explain your answer.}}
  {\ealpara{Nie. $27=3\times 9$, więc ma więcej niż dwa \ealkeytr{dzielniki}.}{No. $27=3\times 9$, so it has more than two \ealkey{factors}.}}

\qq[\LARGE]{\ealpara{Zapisz pierwszą \ealkeytr{liczbę pierwszą} większą od $30$}{Write down the first \ealkey{prime} number greater than $30$}}
  {\ealpara{$31$}{$31$}}

\qq[\LARGE]{\ealpara{Czy $91$ jest \ealkeytr{liczbą pierwszą}?}{Is $91$ a \ealkey{prime} number?}}
  {\ealpara{Nie. $91=7\times 13$.}{No. $91=7\times 13$.}}

\qq[\LARGE]{\ealpara{Ile \ealkeytr{liczb pierwszych} jest między $20$ a $40$?}{How many \ealkey{prime} numbers are there between $20$ and $40$?}}
  {\ealpara{$23,\ 29,\ 31,\ 37$\\[0.3em]Cztery.}{$23,\ 29,\ 31,\ 37$\\[0.3em]Four of them.}}

\qq[\LARGE]{\ealpara{Czy $143$ jest \ealkeytr{liczbą pierwszą}?}{Is $143$ a \ealkey{prime} number?}}
  {\ealpara{Nie. $143=11\times 13$.}{No. $143=11\times 13$.}}

\qq[\LARGE]{\ealpara{Znajdź dwie \ealkeytr{liczby pierwsze}, które sumują się do $24$}{Find two \ealkey{prime} numbers that add up to $24$}}
  {\ealpara{$5+19$ \quad $7+17$ \quad $11+13$}{$5+19$ \quad $7+17$ \quad $11+13$}}

\qq[\Large][\Large]{\ealpara{Aby sprawdzić, czy $211$ jest \ealkeytr{liczbą pierwszą},\\[0.3em]
przez jakie liczby musisz podzielić?\\[0.4em]Czy to \ealkeytr{liczba pierwsza}?}{To test whether $211$ is \ealkey{prime},\\[0.3em]
which numbers do you need to divide by?\\[0.4em]Is it \ealkey{prime}?}}
  {\ealpara{Tylko \ealkeytr{liczby pierwsze} do $\sqrt{211}\approx 14.5$:\\[0.3em]
   $2,\ 3,\ 5,\ 7,\ 11,\ 13$\\[0.3em]
   Żadna z nich nie dzieli $211$, więc $211$ jest \ealkeytr{liczbą pierwszą}.}{Only the \ealkey{primes} up to $\sqrt{211}\approx 14.5$:\\[0.3em]
   $2,\ 3,\ 5,\ 7,\ 11,\ 13$\\[0.3em]
   None of them divides $211$, so $211$ is \ealkey{prime}.}}

\qq[\Large][\Large]{\ealpara{Wyjaśnij, dlaczego $n^2-1$ nigdy nie może być \ealkeytr{liczbą pierwszą}, gdy $n>2$}{Explain why $n^2-1$ can never be \ealkey{prime} when $n>2$}}
  {\ealpara{$n^2-1=(n-1)(n+1)$\\[0.3em]
   Gdy $n>2$, oba nawiasy są większe od $1$,
   więc liczba zawsze ma więcej niż dwa \ealkeytr{dzielniki}.}{$n^2-1=(n-1)(n+1)$\\[0.3em]
   When $n>2$ both brackets are bigger than $1$,
   so the number always has more than two \ealkey{factors}.}}

%================================================================
\section{Is it prime?}
\begin{frame}
  \vfill
  \begin{center}
    \ealpara{Czy to \ealkeytr{liczba pierwsza}?}{Is it \ealkey{prime}?}
    \par\vspace{1em}
    {\usebeamercolor[fg]{structure}\Huge\bfseries Is it \ealkey{prime}?}
  \end{center}
  \vfill
\end{frame}
%================================================================

% Every whole number from 1 to 100, one to a slide, followed by all
% of its factors and the reason it is or is not a prime number.

\prq{1}{$1$}
  {\ealpara{$1$ ma tylko jeden \ealkeytr{dzielnik}. \ealkeytr{Liczba pierwsza} ma dokładnie dwa różne \ealkeytr{dzielniki}, $1$ i samą siebie, a dla $1$ te \ealkeytr{dzielniki} byłyby takie same, więc $1$ nie jest \ealkeytr{liczbą pierwszą}.}{$1$ has only one \ealkey{factor}. A \ealkey{prime} number has exactly two different \ealkey{factors}, $1$ and itself, and for $1$ those would be the same \ealkey{factor}, so $1$ is not a \ealkey{prime} number.}}

\prq{2}{$1$, $2$}
  {\ealpara{$2$ ma dokładnie dwa \ealkeytr{dzielniki}, $1$ i samą siebie, więc $2$ jest \ealkeytr{liczbą pierwszą}. To jedyna parzysta \ealkeytr{liczba pierwsza}, ponieważ każda inna liczba parzysta ma $2$ jako trzeci \ealkeytr{dzielnik}.}{$2$ has exactly two \ealkey{factors}, $1$ and itself, so $2$ is a \ealkey{prime} number. It is the only even \ealkey{prime} number, because every other even number has $2$ as a third \ealkey{factor}.}}

\prq{3}{$1$, $3$}
  {\ealpara{$3$ ma dokładnie dwa \ealkeytr{dzielniki}, $1$ i samą siebie, więc $3$ jest \ealkeytr{liczbą pierwszą}.}{$3$ has exactly two \ealkey{factors}, $1$ and itself, so $3$ is a \ealkey{prime} number.}}

\prq{4}{$1$, $2$, $4$}
  {\ealpara{$4$ ma trzy \ealkeytr{dzielniki}, nie dwa, więc $4$ nie jest \ealkeytr{liczbą pierwszą}.}{$4$ has three \ealkey{factors}, not two, so $4$ is not a \ealkey{prime} number.}}

\prq{5}{$1$, $5$}
  {\ealpara{$5$ ma dokładnie dwa \ealkeytr{dzielniki}, $1$ i samą siebie, więc $5$ jest \ealkeytr{liczbą pierwszą}.\\[0.45em]Musimy sprawdzić tylko $2$ jako \ealkeytr{dzielnik}, ponieważ $3\times 3=9$ jest większe niż $5$.}{$5$ has exactly two \ealkey{factors}, $1$ and itself, so $5$ is a \ealkey{prime} number.\\[0.45em]We only need to test $2$ as a \ealkey{factor}, because $3\times 3=9$ is bigger than $5$.}}

\prq{6}{$1$, $2$, $3$, $6$}
  {\ealpara{$6$ ma cztery \ealkeytr{dzielniki}, nie dwa, więc $6$ nie jest \ealkeytr{liczbą pierwszą}.}{$6$ has four \ealkey{factors}, not two, so $6$ is not a \ealkey{prime} number.}}

\prq{7}{$1$, $7$}
  {\ealpara{$7$ ma dokładnie dwa \ealkeytr{dzielniki}, $1$ i samą siebie, więc $7$ jest \ealkeytr{liczbą pierwszą}.\\[0.45em]Musimy sprawdzić tylko $2$ jako \ealkeytr{dzielnik}, ponieważ $3\times 3=9$ jest większe niż $7$.}{$7$ has exactly two \ealkey{factors}, $1$ and itself, so $7$ is a \ealkey{prime} number.\\[0.45em]We only need to test $2$ as a \ealkey{factor}, because $3\times 3=9$ is bigger than $7$.}}

\prq{8}{$1$, $2$, $4$, $8$}
  {\ealpara{$8$ ma cztery \ealkeytr{dzielniki}, nie dwa, więc $8$ nie jest \ealkeytr{liczbą pierwszą}.}{$8$ has four \ealkey{factors}, not two, so $8$ is not a \ealkey{prime} number.}}

\prq{9}{$1$, $3$, $9$}
  {\ealpara{$9$ ma trzy \ealkeytr{dzielniki}, nie dwa, więc $9$ nie jest \ealkeytr{liczbą pierwszą}.}{$9$ has three \ealkey{factors}, not two, so $9$ is not a \ealkey{prime} number.}}

\prq{10}{$1$, $2$, $5$, $10$}
  {\ealpara{$10$ ma cztery \ealkeytr{dzielniki}, nie dwa, więc $10$ nie jest \ealkeytr{liczbą pierwszą}.}{$10$ has four \ealkey{factors}, not two, so $10$ is not a \ealkey{prime} number.}}

\prq{11}{$1$, $11$}
  {\ealpara{$11$ ma dokładnie dwa \ealkeytr{dzielniki}, $1$ i samą siebie, więc $11$ jest \ealkeytr{liczbą pierwszą}.\\[0.45em]Musimy sprawdzić tylko $2$ i $3$ jako \ealkeytr{dzielniki}, ponieważ $5\times 5=25$ jest większe niż $11$.}{$11$ has exactly two \ealkey{factors}, $1$ and itself, so $11$ is a \ealkey{prime} number.\\[0.45em]We only need to test $2$ and $3$ as \ealkey{factors}, because $5\times 5=25$ is bigger than $11$.}}

\prq{12}{$1$, $2$, $3$, $4$, $6$, $12$}
  {\ealpara{$12$ ma sześć \ealkeytr{dzielników}, nie dwa, więc $12$ nie jest \ealkeytr{liczbą pierwszą}.}{$12$ has six \ealkey{factors}, not two, so $12$ is not a \ealkey{prime} number.}}

\prq{13}{$1$, $13$}
  {\ealpara{$13$ ma dokładnie dwa \ealkeytr{dzielniki}, $1$ i samą siebie, więc $13$ jest \ealkeytr{liczbą pierwszą}.\\[0.45em]Musimy sprawdzić tylko $2$ i $3$ jako \ealkeytr{dzielniki}, ponieważ $5\times 5=25$ jest większe niż $13$.}{$13$ has exactly two \ealkey{factors}, $1$ and itself, so $13$ is a \ealkey{prime} number.\\[0.45em]We only need to test $2$ and $3$ as \ealkey{factors}, because $5\times 5=25$ is bigger than $13$.}}

\prq{14}{$1$, $2$, $7$, $14$}
  {\ealpara{$14$ ma cztery \ealkeytr{dzielniki}, nie dwa, więc $14$ nie jest \ealkeytr{liczbą pierwszą}.}{$14$ has four \ealkey{factors}, not two, so $14$ is not a \ealkey{prime} number.}}

\prq{15}{$1$, $3$, $5$, $15$}
  {\ealpara{$15$ ma cztery \ealkeytr{dzielniki}, nie dwa, więc $15$ nie jest \ealkeytr{liczbą pierwszą}.}{$15$ has four \ealkey{factors}, not two, so $15$ is not a \ealkey{prime} number.}}

\prq{16}{$1$, $2$, $4$, $8$, $16$}
  {\ealpara{$16$ ma pięć \ealkeytr{dzielników}, nie dwa, więc $16$ nie jest \ealkeytr{liczbą pierwszą}.}{$16$ has five \ealkey{factors}, not two, so $16$ is not a \ealkey{prime} number.}}

\prq{17}{$1$, $17$}
  {\ealpara{$17$ ma dokładnie dwa \ealkeytr{dzielniki}, $1$ i samą siebie, więc $17$ jest \ealkeytr{liczbą pierwszą}.\\[0.45em]Musimy sprawdzić tylko $2$ i $3$ jako \ealkeytr{dzielniki}, ponieważ $5\times 5=25$ jest większe niż $17$.}{$17$ has exactly two \ealkey{factors}, $1$ and itself, so $17$ is a \ealkey{prime} number.\\[0.45em]We only need to test $2$ and $3$ as \ealkey{factors}, because $5\times 5=25$ is bigger than $17$.}}

\prq{18}{$1$, $2$, $3$, $6$, $9$, $18$}
  {\ealpara{$18$ ma sześć \ealkeytr{dzielników}, nie dwa, więc $18$ nie jest \ealkeytr{liczbą pierwszą}.}{$18$ has six \ealkey{factors}, not two, so $18$ is not a \ealkey{prime} number.}}

\prq{19}{$1$, $19$}
  {\ealpara{$19$ ma dokładnie dwa \ealkeytr{dzielniki}, $1$ i samą siebie, więc $19$ jest \ealkeytr{liczbą pierwszą}.\\[0.45em]Musimy sprawdzić tylko $2$ i $3$ jako \ealkeytr{dzielniki}, ponieważ $5\times 5=25$ jest większe niż $19$.}{$19$ has exactly two \ealkey{factors}, $1$ and itself, so $19$ is a \ealkey{prime} number.\\[0.45em]We only need to test $2$ and $3$ as \ealkey{factors}, because $5\times 5=25$ is bigger than $19$.}}

\prq{20}{$1$, $2$, $4$, $5$, $10$, $20$}
  {\ealpara{$20$ ma sześć \ealkeytr{dzielników}, nie dwa, więc $20$ nie jest \ealkeytr{liczbą pierwszą}.}{$20$ has six \ealkey{factors}, not two, so $20$ is not a \ealkey{prime} number.}}

\prq{21}{$1$, $3$, $7$, $21$}
  {\ealpara{$21$ ma cztery \ealkeytr{dzielniki}, nie dwa, więc $21$ nie jest \ealkeytr{liczbą pierwszą}.}{$21$ has four \ealkey{factors}, not two, so $21$ is not a \ealkey{prime} number.}}

\prq{22}{$1$, $2$, $11$, $22$}
  {\ealpara{$22$ ma cztery \ealkeytr{dzielniki}, nie dwa, więc $22$ nie jest \ealkeytr{liczbą pierwszą}.}{$22$ has four \ealkey{factors}, not two, so $22$ is not a \ealkey{prime} number.}}

\prq{23}{$1$, $23$}
  {\ealpara{$23$ ma dokładnie dwa \ealkeytr{dzielniki}, $1$ i samą siebie, więc $23$ jest \ealkeytr{liczbą pierwszą}.\\[0.45em]Musimy sprawdzić tylko $2$ i $3$ jako \ealkeytr{dzielniki}, ponieważ $5\times 5=25$ jest większe niż $23$.}{$23$ has exactly two \ealkey{factors}, $1$ and itself, so $23$ is a \ealkey{prime} number.\\[0.45em]We only need to test $2$ and $3$ as \ealkey{factors}, because $5\times 5=25$ is bigger than $23$.}}

\prq{24}{$1$, $2$, $3$, $4$, $6$, $8$, $12$, $24$}
  {\ealpara{$24$ ma osiem \ealkeytr{dzielników}, nie dwa, więc $24$ nie jest \ealkeytr{liczbą pierwszą}.}{$24$ has eight \ealkey{factors}, not two, so $24$ is not a \ealkey{prime} number.}}

\prq{25}{$1$, $5$, $25$}
  {\ealpara{$25$ ma trzy \ealkeytr{dzielniki}, nie dwa, więc $25$ nie jest \ealkeytr{liczbą pierwszą}.}{$25$ has three \ealkey{factors}, not two, so $25$ is not a \ealkey{prime} number.}}

\prq{26}{$1$, $2$, $13$, $26$}
  {\ealpara{$26$ ma cztery \ealkeytr{dzielniki}, nie dwa, więc $26$ nie jest \ealkeytr{liczbą pierwszą}.}{$26$ has four \ealkey{factors}, not two, so $26$ is not a \ealkey{prime} number.}}

\prq{27}{$1$, $3$, $9$, $27$}
  {\ealpara{$27$ ma cztery \ealkeytr{dzielniki}, nie dwa, więc $27$ nie jest \ealkeytr{liczbą pierwszą}.}{$27$ has four \ealkey{factors}, not two, so $27$ is not a \ealkey{prime} number.}}

\prq{28}{$1$, $2$, $4$, $7$, $14$, $28$}
  {\ealpara{$28$ ma sześć \ealkeytr{dzielników}, nie dwa, więc $28$ nie jest \ealkeytr{liczbą pierwszą}.}{$28$ has six \ealkey{factors}, not two, so $28$ is not a \ealkey{prime} number.}}

\prq{29}{$1$, $29$}
  {\ealpara{$29$ ma dokładnie dwa \ealkeytr{dzielniki}, $1$ i samą siebie, więc $29$ jest \ealkeytr{liczbą pierwszą}.\\[0.45em]Musimy sprawdzić tylko $2$, $3$ i $5$ jako \ealkeytr{dzielniki}, ponieważ $7\times 7=49$ jest większe niż $29$.}{$29$ has exactly two \ealkey{factors}, $1$ and itself, so $29$ is a \ealkey{prime} number.\\[0.45em]We only need to test $2$, $3$ and $5$ as \ealkey{factors}, because $7\times 7=49$ is bigger than $29$.}}

\prq{30}{$1$, $2$, $3$, $5$, $6$, $10$, $15$, $30$}
  {\ealpara{$30$ ma osiem \ealkeytr{dzielników}, nie dwa, więc $30$ nie jest \ealkeytr{liczbą pierwszą}.}{$30$ has eight \ealkey{factors}, not two, so $30$ is not a \ealkey{prime} number.}}

\prq{31}{$1$, $31$}
  {\ealpara{$31$ ma dokładnie dwa \ealkeytr{dzielniki}, $1$ i samą siebie, więc $31$ jest \ealkeytr{liczbą pierwszą}.\\[0.45em]Musimy sprawdzić tylko $2$, $3$ i $5$ jako \ealkeytr{dzielniki}, ponieważ $7\times 7=49$ jest większe niż $31$.}{$31$ has exactly two \ealkey{factors}, $1$ and itself, so $31$ is a \ealkey{prime} number.\\[0.45em]We only need to test $2$, $3$ and $5$ as \ealkey{factors}, because $7\times 7=49$ is bigger than $31$.}}

\prq{32}{$1$, $2$, $4$, $8$, $16$, $32$}
  {\ealpara{$32$ ma sześć \ealkeytr{dzielników}, nie dwa, więc $32$ nie jest \ealkeytr{liczbą pierwszą}.}{$32$ has six \ealkey{factors}, not two, so $32$ is not a \ealkey{prime} number.}}

\prq{33}{$1$, $3$, $11$, $33$}
  {\ealpara{$33$ ma cztery \ealkeytr{dzielniki}, nie dwa, więc $33$ nie jest \ealkeytr{liczbą pierwszą}.}{$33$ has four \ealkey{factors}, not two, so $33$ is not a \ealkey{prime} number.}}

\prq{34}{$1$, $2$, $17$, $34$}
  {\ealpara{$34$ ma cztery \ealkeytr{dzielniki}, nie dwa, więc $34$ nie jest \ealkeytr{liczbą pierwszą}.}{$34$ has four \ealkey{factors}, not two, so $34$ is not a \ealkey{prime} number.}}

\prq{35}{$1$, $5$, $7$, $35$}
  {\ealpara{$35$ ma cztery \ealkeytr{dzielniki}, nie dwa, więc $35$ nie jest \ealkeytr{liczbą pierwszą}.}{$35$ has four \ealkey{factors}, not two, so $35$ is not a \ealkey{prime} number.}}

\prq{36}{$1$, $2$, $3$, $4$, $6$, $9$, $12$, $18$, $36$}
  {\ealpara{$36$ ma dziewięć \ealkeytr{dzielników}, nie dwa, więc $36$ nie jest \ealkeytr{liczbą pierwszą}.}{$36$ has nine \ealkey{factors}, not two, so $36$ is not a \ealkey{prime} number.}}

\prq{37}{$1$, $37$}
  {\ealpara{$37$ ma dokładnie dwa \ealkeytr{dzielniki}, $1$ i samą siebie, więc $37$ jest \ealkeytr{liczbą pierwszą}.\\[0.45em]Musimy sprawdzić tylko $2$, $3$ i $5$ jako \ealkeytr{dzielniki}, ponieważ $7\times 7=49$ jest większe niż $37$.}{$37$ has exactly two \ealkey{factors}, $1$ and itself, so $37$ is a \ealkey{prime} number.\\[0.45em]We only need to test $2$, $3$ and $5$ as \ealkey{factors}, because $7\times 7=49$ is bigger than $37$.}}

\prq{38}{$1$, $2$, $19$, $38$}
  {\ealpara{$38$ ma cztery \ealkeytr{dzielniki}, nie dwa, więc $38$ nie jest \ealkeytr{liczbą pierwszą}.}{$38$ has four \ealkey{factors}, not two, so $38$ is not a \ealkey{prime} number.}}

\prq{39}{$1$, $3$, $13$, $39$}
  {\ealpara{$39$ ma cztery \ealkeytr{dzielniki}, nie dwa, więc $39$ nie jest \ealkeytr{liczbą pierwszą}.}{$39$ has four \ealkey{factors}, not two, so $39$ is not a \ealkey{prime} number.}}

\prq{40}{$1$, $2$, $4$, $5$, $8$, $10$, $20$, $40$}
  {\ealpara{$40$ ma osiem \ealkeytr{dzielników}, nie dwa, więc $40$ nie jest \ealkeytr{liczbą pierwszą}.}{$40$ has eight \ealkey{factors}, not two, so $40$ is not a \ealkey{prime} number.}}

\prq{41}{$1$, $41$}
  {\ealpara{$41$ ma dokładnie dwa \ealkeytr{dzielniki}, $1$ i samą siebie, więc $41$ jest \ealkeytr{liczbą pierwszą}.\\[0.45em]Musimy sprawdzić tylko $2$, $3$ i $5$ jako \ealkeytr{dzielniki}, ponieważ $7\times 7=49$ jest większe niż $41$.}{$41$ has exactly two \ealkey{factors}, $1$ and itself, so $41$ is a \ealkey{prime} number.\\[0.45em]We only need to test $2$, $3$ and $5$ as \ealkey{factors}, because $7\times 7=49$ is bigger than $41$.}}

\prq{42}{$1$, $2$, $3$, $6$, $7$, $14$, $21$, $42$}
  {\ealpara{$42$ ma osiem \ealkeytr{dzielników}, nie dwa, więc $42$ nie jest \ealkeytr{liczbą pierwszą}.}{$42$ has eight \ealkey{factors}, not two, so $42$ is not a \ealkey{prime} number.}}

\prq{43}{$1$, $43$}
  {\ealpara{$43$ ma dokładnie dwa \ealkeytr{dzielniki}, $1$ i samą siebie, więc $43$ jest \ealkeytr{liczbą pierwszą}.\\[0.45em]Musimy sprawdzić tylko $2$, $3$ i $5$ jako \ealkeytr{dzielniki}, ponieważ $7\times 7=49$ jest większe niż $43$.}{$43$ has exactly two \ealkey{factors}, $1$ and itself, so $43$ is a \ealkey{prime} number.\\[0.45em]We only need to test $2$, $3$ and $5$ as \ealkey{factors}, because $7\times 7=49$ is bigger than $43$.}}

\prq{44}{$1$, $2$, $4$, $11$, $22$, $44$}
  {\ealpara{$44$ ma sześć \ealkeytr{dzielników}, nie dwa, więc $44$ nie jest \ealkeytr{liczbą pierwszą}.}{$44$ has six \ealkey{factors}, not two, so $44$ is not a \ealkey{prime} number.}}

\prq{45}{$1$, $3$, $5$, $9$, $15$, $45$}
  {\ealpara{$45$ ma sześć \ealkeytr{dzielników}, nie dwa, więc $45$ nie jest \ealkeytr{liczbą pierwszą}.}{$45$ has six \ealkey{factors}, not two, so $45$ is not a \ealkey{prime} number.}}

\prq{46}{$1$, $2$, $23$, $46$}
  {\ealpara{$46$ ma cztery \ealkeytr{dzielniki}, nie dwa, więc $46$ nie jest \ealkeytr{liczbą pierwszą}.}{$46$ has four \ealkey{factors}, not two, so $46$ is not a \ealkey{prime} number.}}

\prq{47}{$1$, $47$}
  {\ealpara{$47$ ma dokładnie dwa \ealkeytr{dzielniki}, $1$ i samą siebie, więc $47$ jest \ealkeytr{liczbą pierwszą}.\\[0.45em]Musimy sprawdzić tylko $2$, $3$ i $5$ jako \ealkeytr{dzielniki}, ponieważ $7\times 7=49$ jest większe niż $47$.}{$47$ has exactly two \ealkey{factors}, $1$ and itself, so $47$ is a \ealkey{prime} number.\\[0.45em]We only need to test $2$, $3$ and $5$ as \ealkey{factors}, because $7\times 7=49$ is bigger than $47$.}}

\prq{48}{$1$, $2$, $3$, $4$, $6$, $8$, $12$, $16$, $24$, $48$}
  {\ealpara{$48$ ma dziesięć \ealkeytr{dzielników}, nie dwa, więc $48$ nie jest \ealkeytr{liczbą pierwszą}.}{$48$ has ten \ealkey{factors}, not two, so $48$ is not a \ealkey{prime} number.}}

\prq{49}{$1$, $7$, $49$}
  {\ealpara{$49$ ma trzy \ealkeytr{dzielniki}, nie dwa, więc $49$ nie jest \ealkeytr{liczbą pierwszą}.}{$49$ has three \ealkey{factors}, not two, so $49$ is not a \ealkey{prime} number.}}

\prq{50}{$1$, $2$, $5$, $10$, $25$, $50$}
  {\ealpara{$50$ ma sześć \ealkeytr{dzielników}, nie dwa, więc $50$ nie jest \ealkeytr{liczbą pierwszą}.}{$50$ has six \ealkey{factors}, not two, so $50$ is not a \ealkey{prime} number.}}

\prq{51}{$1$, $3$, $17$, $51$}
  {\ealpara{$51$ ma cztery \ealkeytr{dzielniki}, nie dwa, więc $51$ nie jest \ealkeytr{liczbą pierwszą}.}{$51$ has four \ealkey{factors}, not two, so $51$ is not a \ealkey{prime} number.}}

\prq{52}{$1$, $2$, $4$, $13$, $26$, $52$}
  {\ealpara{$52$ ma sześć \ealkeytr{dzielników}, nie dwa, więc $52$ nie jest \ealkeytr{liczbą pierwszą}.}{$52$ has six \ealkey{factors}, not two, so $52$ is not a \ealkey{prime} number.}}

\prq{53}{$1$, $53$}
  {\ealpara{$53$ ma dokładnie dwa \ealkeytr{dzielniki}, $1$ i samą siebie, więc $53$ jest \ealkeytr{liczbą pierwszą}.\\[0.45em]Musimy sprawdzić tylko $2$, $3$, $5$ i $7$ jako \ealkeytr{dzielniki}, ponieważ $11\times 11=121$ jest większe niż $53$.}{$53$ has exactly two \ealkey{factors}, $1$ and itself, so $53$ is a \ealkey{prime} number.\\[0.45em]We only need to test $2$, $3$, $5$ and $7$ as \ealkey{factors}, because $11\times 11=121$ is bigger than $53$.}}

\prq{54}{$1$, $2$, $3$, $6$, $9$, $18$, $27$, $54$}
  {\ealpara{$54$ ma osiem \ealkeytr{dzielników}, nie dwa, więc $54$ nie jest \ealkeytr{liczbą pierwszą}.}{$54$ has eight \ealkey{factors}, not two, so $54$ is not a \ealkey{prime} number.}}

\prq{55}{$1$, $5$, $11$, $55$}
  {\ealpara{$55$ ma cztery \ealkeytr{dzielniki}, nie dwa, więc $55$ nie jest \ealkeytr{liczbą pierwszą}.}{$55$ has four \ealkey{factors}, not two, so $55$ is not a \ealkey{prime} number.}}

\prq{56}{$1$, $2$, $4$, $7$, $8$, $14$, $28$, $56$}
  {\ealpara{$56$ ma osiem \ealkeytr{dzielników}, nie dwa, więc $56$ nie jest \ealkeytr{liczbą pierwszą}.}{$56$ has eight \ealkey{factors}, not two, so $56$ is not a \ealkey{prime} number.}}

\prq{57}{$1$, $3$, $19$, $57$}
  {\ealpara{$57$ ma cztery \ealkeytr{dzielniki}, nie dwa, więc $57$ nie jest \ealkeytr{liczbą pierwszą}.}{$57$ has four \ealkey{factors}, not two, so $57$ is not a \ealkey{prime} number.}}

\prq{58}{$1$, $2$, $29$, $58$}
  {\ealpara{$58$ ma cztery \ealkeytr{dzielniki}, nie dwa, więc $58$ nie jest \ealkeytr{liczbą pierwszą}.}{$58$ has four \ealkey{factors}, not two, so $58$ is not a \ealkey{prime} number.}}

\prq{59}{$1$, $59$}
  {\ealpara{$59$ ma dokładnie dwa \ealkeytr{dzielniki}, $1$ i samą siebie, więc $59$ jest \ealkeytr{liczbą pierwszą}.\\[0.45em]Musimy sprawdzić tylko $2$, $3$, $5$ i $7$ jako \ealkeytr{dzielniki}, ponieważ $11\times 11=121$ jest większe niż $59$.}{$59$ has exactly two \ealkey{factors}, $1$ and itself, so $59$ is a \ealkey{prime} number.\\[0.45em]We only need to test $2$, $3$, $5$ and $7$ as \ealkey{factors}, because $11\times 11=121$ is bigger than $59$.}}

\prq{60}{$1$, $2$, $3$, $4$, $5$, $6$, $10$, $12$, $15$, $20$, $30$, $60$}
  {\ealpara{$60$ ma dwanaście \ealkeytr{dzielników}, nie dwa, więc $60$ nie jest \ealkeytr{liczbą pierwszą}.}{$60$ has twelve \ealkey{factors}, not two, so $60$ is not a \ealkey{prime} number.}}

\prq{61}{$1$, $61$}
  {\ealpara{$61$ ma dokładnie dwa \ealkeytr{dzielniki}, $1$ i samą siebie, więc $61$ jest \ealkeytr{liczbą pierwszą}.\\[0.45em]Musimy sprawdzić tylko $2$, $3$, $5$ i $7$ jako \ealkeytr{dzielniki}, ponieważ $11\times 11=121$ jest większe niż $61$.}{$61$ has exactly two \ealkey{factors}, $1$ and itself, so $61$ is a \ealkey{prime} number.\\[0.45em]We only need to test $2$, $3$, $5$ and $7$ as \ealkey{factors}, because $11\times 11=121$ is bigger than $61$.}}

\prq{62}{$1$, $2$, $31$, $62$}
  {\ealpara{$62$ ma cztery \ealkeytr{dzielniki}, nie dwa, więc $62$ nie jest \ealkeytr{liczbą pierwszą}.}{$62$ has four \ealkey{factors}, not two, so $62$ is not a \ealkey{prime} number.}}

\prq{63}{$1$, $3$, $7$, $9$, $21$, $63$}
  {\ealpara{$63$ ma sześć \ealkeytr{dzielników}, nie dwa, więc $63$ nie jest \ealkeytr{liczbą pierwszą}.}{$63$ has six \ealkey{factors}, not two, so $63$ is not a \ealkey{prime} number.}}

\prq{64}{$1$, $2$, $4$, $8$, $16$, $32$, $64$}
  {\ealpara{$64$ ma siedem \ealkeytr{dzielników}, nie dwa, więc $64$ nie jest \ealkeytr{liczbą pierwszą}.}{$64$ has seven \ealkey{factors}, not two, so $64$ is not a \ealkey{prime} number.}}

\prq{65}{$1$, $5$, $13$, $65$}
  {\ealpara{$65$ ma cztery \ealkeytr{dzielniki}, nie dwa, więc $65$ nie jest \ealkeytr{liczbą pierwszą}.}{$65$ has four \ealkey{factors}, not two, so $65$ is not a \ealkey{prime} number.}}

\prq{66}{$1$, $2$, $3$, $6$, $11$, $22$, $33$, $66$}
  {\ealpara{$66$ ma osiem \ealkeytr{dzielników}, nie dwa, więc $66$ nie jest \ealkeytr{liczbą pierwszą}.}{$66$ has eight \ealkey{factors}, not two, so $66$ is not a \ealkey{prime} number.}}

\prq{67}{$1$, $67$}
  {\ealpara{$67$ ma dokładnie dwa \ealkeytr{dzielniki}, $1$ i samą siebie, więc $67$ jest \ealkeytr{liczbą pierwszą}.\\[0.45em]Musimy sprawdzić tylko $2$, $3$, $5$ i $7$ jako \ealkeytr{dzielniki}, ponieważ $11\times 11=121$ jest większe niż $67$.}{$67$ has exactly two \ealkey{factors}, $1$ and itself, so $67$ is a \ealkey{prime} number.\\[0.45em]We only need to test $2$, $3$, $5$ and $7$ as \ealkey{factors}, because $11\times 11=121$ is bigger than $67$.}}

\prq{68}{$1$, $2$, $4$, $17$, $34$, $68$}
  {\ealpara{$68$ ma sześć \ealkeytr{dzielników}, nie dwa, więc $68$ nie jest \ealkeytr{liczbą pierwszą}.}{$68$ has six \ealkey{factors}, not two, so $68$ is not a \ealkey{prime} number.}}

\prq{69}{$1$, $3$, $23$, $69$}
  {\ealpara{$69$ ma cztery \ealkeytr{dzielniki}, nie dwa, więc $69$ nie jest \ealkeytr{liczbą pierwszą}.}{$69$ has four \ealkey{factors}, not two, so $69$ is not a \ealkey{prime} number.}}

\prq{70}{$1$, $2$, $5$, $7$, $10$, $14$, $35$, $70$}
  {\ealpara{$70$ ma osiem \ealkeytr{dzielników}, nie dwa, więc $70$ nie jest \ealkeytr{liczbą pierwszą}.}{$70$ has eight \ealkey{factors}, not two, so $70$ is not a \ealkey{prime} number.}}

\prq{71}{$1$, $71$}
  {\ealpara{$71$ ma dokładnie dwa \ealkeytr{dzielniki}, $1$ i samą siebie, więc $71$ jest \ealkeytr{liczbą pierwszą}.\\[0.45em]Musimy sprawdzić tylko $2$, $3$, $5$ i $7$ jako \ealkeytr{dzielniki}, ponieważ $11\times 11=121$ jest większe niż $71$.}{$71$ has exactly two \ealkey{factors}, $1$ and itself, so $71$ is a \ealkey{prime} number.\\[0.45em]We only need to test $2$, $3$, $5$ and $7$ as \ealkey{factors}, because $11\times 11=121$ is bigger than $71$.}}

\prq{72}{$1$, $2$, $3$, $4$, $6$, $8$, $9$, $12$, $18$, $24$, $36$, $72$}
  {\ealpara{$72$ ma dwanaście \ealkeytr{dzielników}, nie dwa, więc $72$ nie jest \ealkeytr{liczbą pierwszą}.}{$72$ has twelve \ealkey{factors}, not two, so $72$ is not a \ealkey{prime} number.}}

\prq{73}{$1$, $73$}
  {\ealpara{$73$ ma dokładnie dwa \ealkeytr{dzielniki}, $1$ i samą siebie, więc $73$ jest \ealkeytr{liczbą pierwszą}.\\[0.45em]Musimy sprawdzić tylko $2$, $3$, $5$ i $7$ jako \ealkeytr{dzielniki}, ponieważ $11\times 11=121$ jest większe niż $73$.}{$73$ has exactly two \ealkey{factors}, $1$ and itself, so $73$ is a \ealkey{prime} number.\\[0.45em]We only need to test $2$, $3$, $5$ and $7$ as \ealkey{factors}, because $11\times 11=121$ is bigger than $73$.}}

\prq{74}{$1$, $2$, $37$, $74$}
  {\ealpara{$74$ ma cztery \ealkeytr{dzielniki}, nie dwa, więc $74$ nie jest \ealkeytr{liczbą pierwszą}.}{$74$ has four \ealkey{factors}, not two, so $74$ is not a \ealkey{prime} number.}}

\prq{75}{$1$, $3$, $5$, $15$, $25$, $75$}
  {\ealpara{$75$ ma sześć \ealkeytr{dzielników}, nie dwa, więc $75$ nie jest \ealkeytr{liczbą pierwszą}.}{$75$ has six \ealkey{factors}, not two, so $75$ is not a \ealkey{prime} number.}}

\prq{76}{$1$, $2$, $4$, $19$, $38$, $76$}
  {\ealpara{$76$ ma sześć \ealkeytr{dzielników}, nie dwa, więc $76$ nie jest \ealkeytr{liczbą pierwszą}.}{$76$ has six \ealkey{factors}, not two, so $76$ is not a \ealkey{prime} number.}}

\prq{77}{$1$, $7$, $11$, $77$}
  {\ealpara{$77$ ma cztery \ealkeytr{dzielniki}, nie dwa, więc $77$ nie jest \ealkeytr{liczbą pierwszą}.}{$77$ has four \ealkey{factors}, not two, so $77$ is not a \ealkey{prime} number.}}

\prq{78}{$1$, $2$, $3$, $6$, $13$, $26$, $39$, $78$}
  {\ealpara{$78$ ma osiem \ealkeytr{dzielników}, nie dwa, więc $78$ nie jest \ealkeytr{liczbą pierwszą}.}{$78$ has eight \ealkey{factors}, not two, so $78$ is not a \ealkey{prime} number.}}

\prq{79}{$1$, $79$}
  {\ealpara{$79$ ma dokładnie dwa \ealkeytr{dzielniki}, $1$ i samą siebie, więc $79$ jest \ealkeytr{liczbą pierwszą}.\\[0.45em]Musimy sprawdzić tylko $2$, $3$, $5$ i $7$ jako \ealkeytr{dzielniki}, ponieważ $11\times 11=121$ jest większe niż $79$.}{$79$ has exactly two \ealkey{factors}, $1$ and itself, so $79$ is a \ealkey{prime} number.\\[0.45em]We only need to test $2$, $3$, $5$ and $7$ as \ealkey{factors}, because $11\times 11=121$ is bigger than $79$.}}

\prq{80}{$1$, $2$, $4$, $5$, $8$, $10$, $16$, $20$, $40$, $80$}
  {\ealpara{$80$ ma dziesięć \ealkeytr{dzielników}, nie dwa, więc $80$ nie jest \ealkeytr{liczbą pierwszą}.}{$80$ has ten \ealkey{factors}, not two, so $80$ is not a \ealkey{prime} number.}}

\prq{81}{$1$, $3$, $9$, $27$, $81$}
  {\ealpara{$81$ ma pięć \ealkeytr{dzielników}, nie dwa, więc $81$ nie jest \ealkeytr{liczbą pierwszą}.}{$81$ has five \ealkey{factors}, not two, so $81$ is not a \ealkey{prime} number.}}

\prq{82}{$1$, $2$, $41$, $82$}
  {\ealpara{$82$ ma cztery \ealkeytr{dzielniki}, nie dwa, więc $82$ nie jest \ealkeytr{liczbą pierwszą}.}{$82$ has four \ealkey{factors}, not two, so $82$ is not a \ealkey{prime} number.}}

\prq{83}{$1$, $83$}
  {\ealpara{$83$ ma dokładnie dwa \ealkeytr{dzielniki}, $1$ i samą siebie, więc $83$ jest \ealkeytr{liczbą pierwszą}.\\[0.45em]Musimy sprawdzić tylko $2$, $3$, $5$ i $7$ jako \ealkeytr{dzielniki}, ponieważ $11\times 11=121$ jest większe niż $83$.}{$83$ has exactly two \ealkey{factors}, $1$ and itself, so $83$ is a \ealkey{prime} number.\\[0.45em]We only need to test $2$, $3$, $5$ and $7$ as \ealkey{factors}, because $11\times 11=121$ is bigger than $83$.}}

\prq{84}{$1$, $2$, $3$, $4$, $6$, $7$, $12$, $14$, $21$, $28$, $42$, $84$}
  {\ealpara{$84$ ma dwanaście \ealkeytr{dzielników}, nie dwa, więc $84$ nie jest \ealkeytr{liczbą pierwszą}.}{$84$ has twelve \ealkey{factors}, not two, so $84$ is not a \ealkey{prime} number.}}

\prq{85}{$1$, $5$, $17$, $85$}
  {\ealpara{$85$ ma cztery \ealkeytr{dzielniki}, nie dwa, więc $85$ nie jest \ealkeytr{liczbą pierwszą}.}{$85$ has four \ealkey{factors}, not two, so $85$ is not a \ealkey{prime} number.}}

\prq{86}{$1$, $2$, $43$, $86$}
  {\ealpara{$86$ ma cztery \ealkeytr{dzielniki}, nie dwa, więc $86$ nie jest \ealkeytr{liczbą pierwszą}.}{$86$ has four \ealkey{factors}, not two, so $86$ is not a \ealkey{prime} number.}}

\prq{87}{$1$, $3$, $29$, $87$}
  {\ealpara{$87$ ma cztery \ealkeytr{dzielniki}, nie dwa, więc $87$ nie jest \ealkeytr{liczbą pierwszą}.}{$87$ has four \ealkey{factors}, not two, so $87$ is not a \ealkey{prime} number.}}

\prq{88}{$1$, $2$, $4$, $8$, $11$, $22$, $44$, $88$}
  {\ealpara{$88$ ma osiem \ealkeytr{dzielników}, nie dwa, więc $88$ nie jest \ealkeytr{liczbą pierwszą}.}{$88$ has eight \ealkey{factors}, not two, so $88$ is not a \ealkey{prime} number.}}

\prq{89}{$1$, $89$}
  {\ealpara{$89$ ma dokładnie dwa \ealkeytr{dzielniki}, $1$ i samą siebie, więc $89$ jest \ealkeytr{liczbą pierwszą}.\\[0.45em]Musimy sprawdzić tylko $2$, $3$, $5$ i $7$ jako \ealkeytr{dzielniki}, ponieważ $11\times 11=121$ jest większe niż $89$.}{$89$ has exactly two \ealkey{factors}, $1$ and itself, so $89$ is a \ealkey{prime} number.\\[0.45em]We only need to test $2$, $3$, $5$ and $7$ as \ealkey{factors}, because $11\times 11=121$ is bigger than $89$.}}

\prq{90}{$1$, $2$, $3$, $5$, $6$, $9$, $10$, $15$, $18$, $30$, $45$, $90$}
  {\ealpara{$90$ ma dwanaście \ealkeytr{dzielników}, nie dwa, więc $90$ nie jest \ealkeytr{liczbą pierwszą}.}{$90$ has twelve \ealkey{factors}, not two, so $90$ is not a \ealkey{prime} number.}}

\prq{91}{$1$, $7$, $13$, $91$}
  {\ealpara{$91$ ma cztery \ealkeytr{dzielniki}, nie dwa, więc $91$ nie jest \ealkeytr{liczbą pierwszą}.}{$91$ has four \ealkey{factors}, not two, so $91$ is not a \ealkey{prime} number.}}

\prq{92}{$1$, $2$, $4$, $23$, $46$, $92$}
  {\ealpara{$92$ ma sześć \ealkeytr{dzielników}, nie dwa, więc $92$ nie jest \ealkeytr{liczbą pierwszą}.}{$92$ has six \ealkey{factors}, not two, so $92$ is not a \ealkey{prime} number.}}

\prq{93}{$1$, $3$, $31$, $93$}
  {\ealpara{$93$ ma cztery \ealkeytr{dzielniki}, nie dwa, więc $93$ nie jest \ealkeytr{liczbą pierwszą}.}{$93$ has four \ealkey{factors}, not two, so $93$ is not a \ealkey{prime} number.}}

\prq{94}{$1$, $2$, $47$, $94$}
  {\ealpara{$94$ ma cztery \ealkeytr{dzielniki}, nie dwa, więc $94$ nie jest \ealkeytr{liczbą pierwszą}.}{$94$ has four \ealkey{factors}, not two, so $94$ is not a \ealkey{prime} number.}}

\prq{95}{$1$, $5$, $19$, $95$}
  {\ealpara{$95$ ma cztery \ealkeytr{dzielniki}, nie dwa, więc $95$ nie jest \ealkeytr{liczbą pierwszą}.}{$95$ has four \ealkey{factors}, not two, so $95$ is not a \ealkey{prime} number.}}

\prq{96}{$1$, $2$, $3$, $4$, $6$, $8$, $12$, $16$, $24$, $32$, $48$, $96$}
  {\ealpara{$96$ ma dwanaście \ealkeytr{dzielników}, nie dwa, więc $96$ nie jest \ealkeytr{liczbą pierwszą}.}{$96$ has twelve \ealkey{factors}, not two, so $96$ is not a \ealkey{prime} number.}}

\prq{97}{$1$, $97$}
  {\ealpara{$97$ ma dokładnie dwa \ealkeytr{dzielniki}, $1$ i samą siebie, więc $97$ jest \ealkeytr{liczbą pierwszą}.\\[0.45em]Musimy sprawdzić tylko $2$, $3$, $5$ i $7$ jako \ealkeytr{dzielniki}, ponieważ $11\times 11=121$ jest większe niż $97$.}{$97$ has exactly two \ealkey{factors}, $1$ and itself, so $97$ is a \ealkey{prime} number.\\[0.45em]We only need to test $2$, $3$, $5$ and $7$ as \ealkey{factors}, because $11\times 11=121$ is bigger than $97$.}}

\prq{98}{$1$, $2$, $7$, $14$, $49$, $98$}
  {\ealpara{$98$ ma sześć \ealkeytr{dzielników}, nie dwa, więc $98$ nie jest \ealkeytr{liczbą pierwszą}.}{$98$ has six \ealkey{factors}, not two, so $98$ is not a \ealkey{prime} number.}}

\prq{99}{$1$, $3$, $9$, $11$, $33$, $99$}
  {\ealpara{$99$ ma sześć \ealkeytr{dzielników}, nie dwa, więc $99$ nie jest \ealkeytr{liczbą pierwszą}.}{$99$ has six \ealkey{factors}, not two, so $99$ is not a \ealkey{prime} number.}}

\prq{100}{$1$, $2$, $4$, $5$, $10$, $20$, $25$, $50$, $100$}
  {\ealpara{$100$ ma dziewięć \ealkeytr{dzielników}, nie dwa, więc $100$ nie jest \ealkeytr{liczbą pierwszą}.}{$100$ has nine \ealkey{factors}, not two, so $100$ is not a \ealkey{prime} number.}}

%================================================================
\section{Prime factor decomposition}
\begin{frame}
  \vfill
  \begin{center}
    \ealpara{\ealkeytr{Rozkład na czynniki pierwsze}}{\ealkey{Prime factor decomposition}}
    \par\vspace{1em}
    {\usebeamercolor[fg]{structure}\Huge\bfseries \ealkey{Prime factor decomposition}}
  \end{center}
  \vfill
\end{frame}
%================================================================

\tf[\Large]{\ealpara{$12=2^2\times 3$ to \ealkeytr{rozkład na czynniki pierwsze} $12$}{$12=2^2\times 3$ is the \ealkey{prime factor decomposition} of $12$}}
  {\ealpara{\textbf{Prawda.} $2$ i $3$ są obie \ealkeytr{liczbami pierwszymi},
   a $4\times 3=12$.}{\textbf{True.} $2$ and $3$ are both \ealkey{prime},
   and $4\times 3=12$.}}

\tf[\Large]{\ealpara{$2\times 3\times 4$ to \ealkeytr{rozkład na czynniki pierwsze} $24$}{$2\times 3\times 4$ is the \ealkey{prime factor decomposition} of $24$}}
  {\ealpara{\textbf{Fałsz.} $4$ nie jest \ealkeytr{liczbą pierwszą}.\\[0.3em]
   $24=2^3\times 3$}{\textbf{False.} $4$ is not \ealkey{prime}.\\[0.3em]
   $24=2^3\times 3$}}

\tf[\Large]{\ealpara{Każdą liczbę całkowitą większą od $1$ można zapisać jako \ealkeytr{iloczyn}
\ealkeytr{liczb pierwszych}}{Every whole number greater than $1$ can be written as a \ealkey{product}
of \ealkey{primes}}}
  {\ealpara{\textbf{Prawda.} \ealkeytr{Liczba pierwsza} jest swoim własnym rozkładem, a każda inna liczba
   dzieli się dalej, aż pozostaną tylko \ealkeytr{liczby pierwsze}.}{\textbf{True.} A \ealkey{prime} is its own decomposition, and any other number
   keeps splitting until only \ealkey{primes} are left.}}

\tf[\Large]{\ealpara{Liczba może mieć dwa różne \ealkeytr{rozkłady na czynniki pierwsze}}{A number can have two different \ealkey{prime factor decompositions}}}
  {\ealpara{\textbf{Fałsz.} Poza kolejnością zapisu,
   każda liczba ma dokładnie jeden rozkład.}{\textbf{False.} Apart from the order you write them in,
   every number has exactly one decomposition.}}

\tf[\Large]{\ealpara{$2^3\times 5=30$}{$2^3\times 5=30$}}
  {\ealpara{\textbf{Fałsz.} $2^3\times 5=8\times 5=40$.\\[0.3em]
   $30=2\times 3\times 5$}{\textbf{False.} $2^3\times 5=8\times 5=40$.\\[0.3em]
   $30=2\times 3\times 5$}}

\qq[\LARGE]{\ealpara{Zapisz $12$ jako \ealkeytr{iloczyn} jej \ealkeytr{czynników pierwszych}}{Write $12$ as a \ealkey{product} of its \ealkey{prime factors}}}
  {\ealpara{$12=2\times 2\times 3=2^2\times 3$}{$12=2\times 2\times 3=2^2\times 3$}}

\qq[\LARGE]{\ealpara{Zapisz $18$ jako \ealkeytr{iloczyn} jej \ealkeytr{czynników pierwszych}}{Write $18$ as a \ealkey{product} of its \ealkey{prime factors}}}
  {\ealpara{$18=2\times 3\times 3=2\times 3^2$}{$18=2\times 3\times 3=2\times 3^2$}}

\qq[\LARGE]{\ealpara{Zapisz $60$ jako \ealkeytr{iloczyn} jej \ealkeytr{czynników pierwszych}}{Write $60$ as a \ealkey{product} of its \ealkey{prime factors}}}
  {\ealpara{$60=2\times 2\times 3\times 5=2^2\times 3\times 5$}{$60=2\times 2\times 3\times 5=2^2\times 3\times 5$}}

\qq[\LARGE]{\ealpara{Zapisz $84$ jako \ealkeytr{iloczyn} jej \ealkeytr{czynników pierwszych}}{Write $84$ as a \ealkey{product} of its \ealkey{prime factors}}}
  {\ealpara{$84=2\times 2\times 3\times 7=2^2\times 3\times 7$}{$84=2\times 2\times 3\times 7=2^2\times 3\times 7$}}

\qq[\LARGE]{\ealpara{Zapisz $200$ jako \ealkeytr{iloczyn} jej \ealkeytr{czynników pierwszych}}{Write $200$ as a \ealkey{product} of its \ealkey{prime factors}}}
  {\ealpara{$200=2\times 2\times 2\times 5\times 5=2^3\times 5^2$}{$200=2\times 2\times 2\times 5\times 5=2^3\times 5^2$}}

\qq[\LARGE]{\ealpara{Zapisz $360$ jako \ealkeytr{iloczyn} jej \ealkeytr{czynników pierwszych}}{Write $360$ as a \ealkey{product} of its \ealkey{prime factors}}}
  {\ealpara{$360=2^3\times 3^2\times 5$}{$360=2^3\times 3^2\times 5$}}

\qq[\LARGE]{\ealpara{Zapisz $1001$ jako \ealkeytr{iloczyn} jej \ealkeytr{czynników pierwszych}}{Write $1001$ as a \ealkey{product} of its \ealkey{prime factors}}}
  {\ealpara{$1001=7\times 143=7\times 11\times 13$}{$1001=7\times 143=7\times 11\times 13$}}

\qq[\LARGE]{\ealpara{Oblicz $2^3\times 3\times 5^2$ jako zwykłą liczbę}{Work out $2^3\times 3\times 5^2$ as an ordinary number}}
  {\ealpara{$8\times 3\times 25=600$}{$8\times 3\times 25=600$}}

\qq[\Large][\Large]{\ealpara{$n=2^3\times 3\times 5$.\\[0.4em]
Zapisz $2n$ jako \ealkeytr{iloczyn} \ealkeytr{czynników pierwszych}.}{$n=2^3\times 3\times 5$.\\[0.4em]
Write $2n$ as a \ealkey{product} of \ealkey{prime factors}.}}
  {\ealpara{Podwojenie dodaje jeszcze jeden \ealkeytr{czynnik} $2$:\\[0.3em]
   $2n=2^4\times 3\times 5$}{Doubling adds one more \ealkey{factor} of $2$:\\[0.3em]
   $2n=2^4\times 3\times 5$}}

\qq[\Large][\Large]{\ealpara{$720=2^4\times 3^2\times 5$.\\[0.4em]
Ile \ealkeytr{dzielników} ma $720$?}{$720=2^4\times 3^2\times 5$.\\[0.4em]
How many \ealkey{factors} does $720$ have?}}
  {\ealpara{Każdy \ealkeytr{dzielnik} używa od $0$ do $4$ dwójek, od $0$ do $2$ trójek i $0$ lub $1$ piątki:\\[0.3em]
   $5\times 3\times 2=30$ \ealkeytr{dzielników}}{Each \ealkey{factor} uses $0$ to $4$ twos, $0$ to $2$ threes and $0$ or $1$ five:\\[0.3em]
   $5\times 3\times 2=30$ \ealkey{factors}}}

%================================================================
\section{HCF and LCM from lists}
\begin{frame}
  \vfill
  \begin{center}
    \ealpara{\ealkeytr{Największy wspólny dzielnik} i \ealkeytr{najmniejsza wspólna wielokrotność} z list}{\ealkey{HCF} and \ealkey{LCM} from lists}
    \par\vspace{1em}
    {\usebeamercolor[fg]{structure}\Huge\bfseries \ealkey{HCF} and \ealkey{LCM} from lists}
  \end{center}
  \vfill
\end{frame}
%================================================================

\hcflcm{4}{6}
  {$1$, \hcfmark{2}, $4$}
  {$1$, \hcfmark{2}, $3$, $6$}
  {$4$, $8$, \lcmmark{12}, $16$, $20$,\\[0.45em]
   $24$, $28$, $32$, $36$, $40$}
  {$6$, \lcmmark{12}, $18$, $24$, $30$,\\[0.45em]
   $36$, $42$, $48$, $54$, $60$}

\hcflcm{6}{8}
  {$1$, \hcfmark{2}, $3$, $6$}
  {$1$, \hcfmark{2}, $4$, $8$}
  {$6$, $12$, $18$, \lcmmark{24}, $30$,\\[0.45em]
   $36$, $42$, $48$, $54$, $60$}
  {$8$, $16$, \lcmmark{24}, $32$, $40$,\\[0.45em]
   $48$, $56$, $64$, $72$, $80$}

\hcflcm{6}{9}
  {$1$, $2$, \hcfmark{3}, $6$}
  {$1$, \hcfmark{3}, $9$}
  {$6$, $12$, \lcmmark{18}, $24$, $30$,\\[0.45em]
   $36$, $42$, $48$, $54$, $60$}
  {$9$, \lcmmark{18}, $27$, $36$, $45$,\\[0.45em]
   $54$, $63$, $72$, $81$, $90$}

\hcflcm{8}{12}
  {$1$, $2$, \hcfmark{4}, $8$}
  {$1$, $2$, $3$, \hcfmark{4}, $6$, $12$}
  {$8$, $16$, \lcmmark{24}, $32$, $40$,\\[0.45em]
   $48$, $56$, $64$, $72$, $80$}
  {$12$, \lcmmark{24}, $36$, $48$, $60$,\\[0.45em]
   $72$, $84$, $96$, $108$, $120$}

\hcflcm{10}{15}
  {$1$, $2$, \hcfmark{5}, $10$}
  {$1$, $3$, \hcfmark{5}, $15$}
  {$10$, $20$, \lcmmark{30}, $40$, $50$,\\[0.45em]
   $60$, $70$, $80$, $90$, $100$}
  {$15$, \lcmmark{30}, $45$, $60$, $75$,\\[0.45em]
   $90$, $105$, $120$, $135$, $150$}

\hcflcm{9}{12}
  {$1$, \hcfmark{3}, $9$}
  {$1$, $2$, \hcfmark{3}, $4$, $6$, $12$}
  {$9$, $18$, $27$, \lcmmark{36}, $45$,\\[0.45em]
   $54$, $63$, $72$, $81$, $90$}
  {$12$, $24$, \lcmmark{36}, $48$, $60$,\\[0.45em]
   $72$, $84$, $96$, $108$, $120$}

\hcflcm{12}{18}
  {$1$, $2$, $3$, $4$, \hcfmark{6}, $12$}
  {$1$, $2$, $3$, \hcfmark{6}, $9$, $18$}
  {$12$, $24$, \lcmmark{36}, $48$, $60$,\\[0.45em]
   $72$, $84$, $96$, $108$, $120$}
  {$18$, \lcmmark{36}, $54$, $72$, $90$,\\[0.45em]
   $108$, $126$, $144$, $162$, $180$}

\hcflcm{14}{21}
  {$1$, $2$, \hcfmark{7}, $14$}
  {$1$, $3$, \hcfmark{7}, $21$}
  {$14$, $28$, \lcmmark{42}, $56$, $70$,\\[0.45em]
   $84$, $98$, $112$, $126$, $140$}
  {$21$, \lcmmark{42}, $63$, $84$, $105$,\\[0.45em]
   $126$, $147$, $168$, $189$, $210$}

\hcflcm{15}{20}
  {$1$, $3$, \hcfmark{5}, $15$}
  {$1$, $2$, $4$, \hcfmark{5}, $10$, $20$}
  {$15$, $30$, $45$, \lcmmark{60}, $75$,\\[0.45em]
   $90$, $105$, $120$, $135$, $150$}
  {$20$, $40$, \lcmmark{60}, $80$, $100$,\\[0.45em]
   $120$, $140$, $160$, $180$, $200$}

\hcflcm{8}{9}
  {\hcfmark{1}, $2$, $4$, $8$}
  {\hcfmark{1}, $3$, $9$}
  {$8$, $16$, $24$, $32$, $40$,\\[0.45em]
   $48$, $56$, $64$, \lcmmark{72}, $80$}
  {$9$, $18$, $27$, $36$, $45$,\\[0.45em]
   $54$, $63$, \lcmmark{72}, $81$, $90$}

%================================================================
\section{HCF and LCM by prime factor decomposition}
\begin{frame}
  \vfill
  \begin{center}
    \ealpara{\ealkeytr{Największy wspólny dzielnik} i \ealkeytr{najmniejsza wspólna wielokrotność} przez \ealkeytr{rozkład na czynniki pierwsze}}{\ealkey{HCF} and \ealkey{LCM} by \ealkey{prime factor decomposition}}
    \par\vspace{1em}
    {\usebeamercolor[fg]{structure}\Huge\bfseries \ealkey{HCF} and \ealkey{LCM} by \ealkey{prime factor decomposition}}
  \end{center}
  \vfill
\end{frame}
%================================================================

\tf[\Large]{\ealpara{\ealkeytr{Największy wspólny dzielnik} dwóch liczb jest zawsze mniejszy od obu z nich}{\ealkey{The HCF} of two numbers is always smaller than both of them}}
  {\ealpara{\textbf{Fałsz.} \ealkeytr{HCF} $6$ i $12$ to $6$,
   która jest jedną z tych liczb.}{\textbf{False.} The \ealkey{HCF} of $6$ and $12$ is $6$,
   which is one of the numbers itself.}}

\tf[\Large]{\ealpara{\ealkeytr{Najmniejsza wspólna wielokrotność} dwóch liczb jest zawsze \ealkeytr{wielokrotnością} obu z nich}{\ealkey{The LCM} of two numbers is always a \ealkey{multiple} of both of them}}
  {\ealpara{\textbf{Prawda.} To jest dokładnie definicja wspólnej \ealkeytr{wielokrotności};
   \ealkeytr{LCM} jest najmniejszą z nich.}{\textbf{True.} That is exactly what a common \ealkey{multiple} is;
   the \ealkey{LCM} is the smallest one.}}

\tf[\Large]{\ealpara{Jeśli $a=2^2\times 3$ i $b=2\times 3^2$,\\[0.3em]to \ealkeytr{HCF}
$a$ i $b$ wynosi $2\times 3$}{If $a=2^2\times 3$ and $b=2\times 3^2$,\\[0.3em]then the \ealkey{HCF} of
$a$ and $b$ is $2\times 3$}}
  {\ealpara{\textbf{Prawda.} Bierzemy \emph{najniższą} \ealkeytr{potęgę} każdej \ealkeytr{liczby pierwszej}, która występuje
   w obu: $2^1\times 3^1=6$.}{\textbf{True.} Take the \emph{lowest} \ealkey{power} of each \ealkey{prime} that appears
   in both: $2^1\times 3^1=6$.}}

\tf[\Large]{\ealpara{$\mathrm{HCF}\times\mathrm{LCM}$ daje \ealkeytr{iloczyn} tych dwóch liczb}{$\mathrm{HCF}\times\mathrm{LCM}$ gives the \ealkey{product} of the two numbers}}
  {\ealpara{\textbf{Prawda.} Każda \ealkeytr{liczba pierwsza} jest liczona raz w najniższej \ealkeytr{potędze} i raz
   w najwyższej, co jest tym samym, co policzenie obu liczb.}{\textbf{True.} Every \ealkey{prime} is counted once at its lowest \ealkey{power} and once
   at its highest, which is the same as counting both numbers.}}

\tf[\Large]{\ealpara{\ealkeytr{Najmniejsza wspólna wielokrotność} $6$ i $8$ wynosi $48$}{\ealkey{The lowest common multiple} of $6$ and $8$ is $48$}}
  {\ealpara{\textbf{Fałsz.} $6=2\times 3$ i $8=2^3$,
   więc \ealkeytr{LCM} to $2^3\times 3=24$.}{\textbf{False.} $6=2\times 3$ and $8=2^3$,
   so the \ealkey{LCM} is $2^3\times 3=24$.}}

\qq[\Large][\Large]{\ealpara{$12=2^2\times 3$ \quad i \quad $18=2\times 3^2$\\[0.4em]
Znajdź \ealkeytr{HCF} $12$ i $18$.}{$12=2^2\times 3$ \quad and \quad $18=2\times 3^2$\\[0.4em]
Find the \ealkey{HCF} of $12$ and $18$.}}
  {\ealpara{Najniższa \ealkeytr{potęga} każdej wspólnej \ealkeytr{liczby pierwszej}:\\[0.3em]
   $\mathrm{HCF}=2\times 3=6$}{Lowest \ealkey{power} of each shared \ealkey{prime}:\\[0.3em]
   $\mathrm{HCF}=2\times 3=6$}}

\qq[\Large][\Large]{\ealpara{$12=2^2\times 3$ \quad i \quad $18=2\times 3^2$\\[0.4em]
Znajdź \ealkeytr{LCM} $12$ i $18$.}{$12=2^2\times 3$ \quad and \quad $18=2\times 3^2$\\[0.4em]
Find the \ealkey{LCM} of $12$ and $18$.}}
  {\ealpara{Najwyższa \ealkeytr{potęga} każdej \ealkeytr{liczby pierwszej}:\\[0.3em]
   $\mathrm{LCM}=2^2\times 3^2=36$}{Highest \ealkey{power} of each \ealkey{prime}:\\[0.3em]
   $\mathrm{LCM}=2^2\times 3^2=36$}}

\qq[\Large][\Large]{\ealpara{$24=2^3\times 3$ \quad i \quad $60=2^2\times 3\times 5$\\[0.4em]
Znajdź \ealkeytr{HCF} $24$ i $60$.}{$24=2^3\times 3$ \quad and \quad $60=2^2\times 3\times 5$\\[0.4em]
Find the \ealkey{HCF} of $24$ and $60$.}}
  {\ealpara{$\mathrm{HCF}=2^2\times 3=12$}{$\mathrm{HCF}=2^2\times 3=12$}}

\qq[\Large][\Large]{\ealpara{$24=2^3\times 3$ \quad i \quad $60=2^2\times 3\times 5$\\[0.4em]
Znajdź \ealkeytr{LCM} $24$ i $60$.}{$24=2^3\times 3$ \quad and \quad $60=2^2\times 3\times 5$\\[0.4em]
Find the \ealkey{LCM} of $24$ and $60$.}}
  {\ealpara{$\mathrm{LCM}=2^3\times 3\times 5=120$}{$\mathrm{LCM}=2^3\times 3\times 5=120$}}

\qq[\Large][\large]{\ealpara{Znajdź \ealkeytr{HCF} i \ealkeytr{LCM} $84$ i $120$}{Find the \ealkey{HCF} and the \ealkey{LCM} of $84$ and $120$}}
  {\ealpara{$84=2^2\times 3\times 7$ \qquad $120=2^3\times 3\times 5$\\[0.4em]
   $\mathrm{HCF}=2^2\times 3=12$\\[0.3em]
   $\mathrm{LCM}=2^3\times 3\times 5\times 7=840$}{$84=2^2\times 3\times 7$ \qquad $120=2^3\times 3\times 5$\\[0.4em]
   $\mathrm{HCF}=2^2\times 3=12$\\[0.3em]
   $\mathrm{LCM}=2^3\times 3\times 5\times 7=840$}}

\qq[\Large][\large]{\ealpara{Znajdź \ealkeytr{HCF} i \ealkeytr{LCM} $126$ i $180$}{Find the \ealkey{HCF} and the \ealkey{LCM} of $126$ and $180$}}
  {\ealpara{$126=2\times 3^2\times 7$ \qquad $180=2^2\times 3^2\times 5$\\[0.4em]
   $\mathrm{HCF}=2\times 3^2=18$\\[0.3em]
   $\mathrm{LCM}=2^2\times 3^2\times 5\times 7=1260$}{$126=2\times 3^2\times 7$ \qquad $180=2^2\times 3^2\times 5$\\[0.4em]
   $\mathrm{HCF}=2\times 3^2=18$\\[0.3em]
   $\mathrm{LCM}=2^2\times 3^2\times 5\times 7=1260$}}

\qq[\Large][\large]{\ealpara{$a=2^3\times 3^2\times 5$ \quad i \quad
$b=2^2\times 3\times 5^2$\\[0.4em]
Znajdź \ealkeytr{HCF} i \ealkeytr{LCM}, w \ealkeytr{postaci potęgowej}.}{$a=2^3\times 3^2\times 5$ \quad and \quad
$b=2^2\times 3\times 5^2$\\[0.4em]
Find the \ealkey{HCF} and the \ealkey{LCM}, in \ealkey{index form}.}}
  {\ealpara{$\mathrm{HCF}=2^2\times 3\times 5$\\[0.3em]
   $\mathrm{LCM}=2^3\times 3^2\times 5^2$}{$\mathrm{HCF}=2^2\times 3\times 5$\\[0.3em]
   $\mathrm{LCM}=2^3\times 3^2\times 5^2$}}

\qq[\Large][\large]{\ealpara{Dwie liczby mają \ealkeytr{HCF} $6$ i \ealkeytr{LCM} $72$.\\[0.4em]
Jedna z tych liczb to $24$. Znajdź drugą.}{Two numbers have \ealkey{HCF} $6$ and \ealkey{LCM} $72$.\\[0.4em]
One of the numbers is $24$. Find the other.}}
  {\ealpara{$\mathrm{HCF}\times\mathrm{LCM}=6\times 72=432$, co jest
   \ealkeytr{iloczynem} tych dwóch liczb.\\[0.4em]
   $432\div 24=18$}{$\mathrm{HCF}\times\mathrm{LCM}=6\times 72=432$, which is the
   \ealkey{product} of the two numbers.\\[0.4em]
   $432\div 24=18$}}

\qq[\Large][\large]{\ealpara{Znajdź \ealkeytr{HCF} i \ealkeytr{LCM} z\\[0.4em]
$2^5\times 3^2\times 11$ \quad i \quad $2^3\times 3^4\times 7$\\[0.4em]
Zostaw odpowiedzi w \ealkeytr{postaci potęgowej}.}{Find the \ealkey{HCF} and the \ealkey{LCM} of\\[0.4em]
$2^5\times 3^2\times 11$ \quad and \quad $2^3\times 3^4\times 7$\\[0.4em]
Leave your answers in \ealkey{index form}.}}
  {\ealpara{$\mathrm{HCF}=2^3\times 3^2$\\[0.3em]
   $\mathrm{LCM}=2^5\times 3^4\times 7\times 11$}{$\mathrm{HCF}=2^3\times 3^2$\\[0.3em]
   $\mathrm{LCM}=2^5\times 3^4\times 7\times 11$}}

\qq[\Large][\large]{\ealpara{Dwie liczby mają \ealkeytr{HCF} $12$ i \ealkeytr{LCM} $180$.\\[0.4em]
Znajdź te liczby. Czy jest więcej niż jedna odpowiedź?}{Two numbers have \ealkey{HCF} $12$ and \ealkey{LCM} $180$.\\[0.4em]
Find the numbers. Is there more than one answer?}}
  {\ealpara{Ich \ealkeytr{iloczyn} to $12\times 180=2160$, a obie są \ealkeytr{wielokrotnościami} $12$.\\[0.4em]
   $12$ i $180$ \quad lub \quad $36$ i $60$\\[0.3em]
   Obie pary pasują, więc jest więcej niż jedna odpowiedź.}{Their \ealkey{product} is $12\times 180=2160$, and both are \ealkey{multiples} of $12$.\\[0.4em]
   $12$ and $180$ \quad or \quad $36$ and $60$\\[0.3em]
   Both pairs work, so there is more than one answer.}}

%================================================================
\section{Square numbers from prime factors}
\begin{frame}
  \vfill
  \begin{center}
    \ealpara{\ealkeytr{Liczby kwadratowe} z \ealkeytr{czynników pierwszych}}{\ealkey{Square numbers} from \ealkey{prime factors}}
    \par\vspace{1em}
    {\usebeamercolor[fg]{structure}\Huge\bfseries \ealkey{Square numbers} from \ealkey{prime factors}}
  \end{center}
  \vfill
\end{frame}
%================================================================

\tf[\Large]{\ealpara{Liczba jest kwadratowa dokładnie wtedy, gdy każda \ealkeytr{potęga} w jej
\ealkeytr{rozkładzie na czynniki pierwsze} jest parzysta}{A number is square exactly when every \ealkey{power} in its \ealkey{prime}
\ealkey{factor decomposition} is even}}
  {\ealpara{\textbf{Prawda.} Połowienie każdej parzystej \ealkeytr{potęgi} daje \ealkeytr{pierwiastek}.}{\textbf{True.} Halving every even \ealkey{power} gives the \ealkey{square root}.}}

\tf[\Large]{\ealpara{$2^3\times 3^2$ jest \ealkeytr{liczbą kwadratową}}{$2^3\times 3^2$ is a \ealkey{square number}}}
  {\ealpara{\textbf{Fałsz.} \ealkeytr{Potęga} $2$ jest nieparzysta,
   więc liczby nie można podzielić na dwie równe połowy.}{\textbf{False.} The \ealkey{power} of $2$ is odd,
   so the number cannot be split into two equal halves.}}

\tf[\Large]{\ealpara{$2^4\times 5^6$ jest \ealkeytr{liczbą kwadratową}}{$2^4\times 5^6$ is a \ealkey{square number}}}
  {\ealpara{\textbf{Prawda.} $2^4\times 5^6=\left(2^2\times 5^3\right)^2$}{\textbf{True.} $2^4\times 5^6=\left(2^2\times 5^3\right)^2$}}

\tf[\Large]{\ealpara{Każda \ealkeytr{liczba kwadratowa} ma parzystą liczbę \ealkeytr{dzielników}}{Every \ealkey{square number} has an even number of \ealkey{factors}}}
  {\ealpara{\textbf{Fałsz.} \ealkeytr{Liczby kwadratowe} mają \emph{nieparzystą} liczbę \ealkeytr{dzielników}:\\[0.3em]
   $16$ ma $1,\ 2,\ 4,\ 8,\ 16$.}{\textbf{False.} \ealkey{Square numbers} have an \emph{odd} number of \ealkey{factors}:\\[0.3em]
   $16$ has $1,\ 2,\ 4,\ 8,\ 16$.}}

\tf[\Large]{\ealpara{Jeśli $n$ jest \ealkeytr{liczbą kwadratową}, to $4n$ też jest \ealkeytr{liczbą kwadratową}}{If $n$ is a \ealkey{square number}, then $4n$ is a \ealkey{square number} too}}
  {\ealpara{\textbf{Prawda.} Mnożenie przez $4=2^2$ utrzymuje każdą \ealkeytr{potęgę} parzystą.}{\textbf{True.} Multiplying by $4=2^2$ keeps every \ealkey{power} even.}}

\qq[\Large][\Large]{\ealpara{Użyj \ealkeytr{czynników pierwszych}, aby pokazać, że $36$ jest \ealkeytr{liczbą kwadratową}}{Use \ealkey{prime factors} to show that $36$ is a \ealkey{square number}}}
  {\ealpara{$36=2^2\times 3^2$\\[0.3em]
   Obie \ealkeytr{potęgi} są parzyste, a $36=(2\times 3)^2=6^2$.}{$36=2^2\times 3^2$\\[0.3em]
   Both \ealkey{powers} are even, and $36=(2\times 3)^2=6^2$.}}

\qq[\Large][\Large]{\ealpara{Czy $2^2\times 3^2\times 5^2$ jest \ealkeytr{liczbą kwadratową}?\\[0.4em]
Jeśli tak, to jaki jest jej \ealkeytr{pierwiastek}?}{Is $2^2\times 3^2\times 5^2$ a \ealkey{square number}?\\[0.4em]
If so, what is its \ealkey{square root}?}}
  {\ealpara{Tak --- każda \ealkeytr{potęga} jest parzysta.\\[0.3em]
   $\sqrt{2^2\times 3^2\times 5^2}=2\times 3\times 5=30$}{Yes --- every \ealkey{power} is even.\\[0.3em]
   $\sqrt{2^2\times 3^2\times 5^2}=2\times 3\times 5=30$}}

\qq[\Large][\Large]{\ealpara{Czy $2^3\times 3^2$ jest \ealkeytr{liczbą kwadratową}?}{Is $2^3\times 3^2$ a \ealkey{square number}?}}
  {\ealpara{Nie. Połowienie \ealkeytr{potęg} wymagałoby $2^{1.5}$,\\[0.3em]
   więc \ealkeytr{potęga} $2$ musi być parzysta, a nie jest.}{No. Halving the \ealkey{powers} would need $2^{1.5}$,\\[0.3em]
   so the \ealkey{power} of $2$ must be even and it is not.}}

\qq[\Large][\Large]{\ealpara{Oblicz $\sqrt{2^6\times 3^4}$}{Work out $\sqrt{2^6\times 3^4}$}}
  {\ealpara{Połów każdą \ealkeytr{potęgę}:\\[0.3em]
   $2^3\times 3^2=8\times 9=72$}{Halve each \ealkey{power}:\\[0.3em]
   $2^3\times 3^2=8\times 9=72$}}

\qq[\Large][\Large]{\ealpara{Użyj \ealkeytr{czynników pierwszych}, aby zdecydować, czy $400$ jest \ealkeytr{liczbą kwadratową}}{Use \ealkey{prime factors} to decide whether $400$ is a \ealkey{square number}}}
  {\ealpara{$400=2^4\times 5^2$\\[0.3em]
   Obie \ealkeytr{potęgi} są parzyste, więc $400=\left(2^2\times 5\right)^2=20^2$.}{$400=2^4\times 5^2$\\[0.3em]
   Both \ealkey{powers} are even, so $400=\left(2^2\times 5\right)^2=20^2$.}}

\qq[\Large][\Large]{\ealpara{Użyj \ealkeytr{czynników pierwszych}, aby zdecydować, czy $150$ jest \ealkeytr{liczbą kwadratową}}{Use \ealkey{prime factors} to decide whether $150$ is a \ealkey{square number}}}
  {\ealpara{$150=2\times 3\times 5^2$\\[0.3em]
   \ealkeytr{Potęgi} $2$ i $3$ są nieparzyste, więc $150$ nie jest kwadratowa.}{$150=2\times 3\times 5^2$\\[0.3em]
   The \ealkey{powers} of $2$ and $3$ are odd, so $150$ is not square.}}

\qq[\Large][\Large]{\ealpara{Oblicz $\sqrt{2^8\times 3^2\times 5^4}$}{Work out $\sqrt{2^8\times 3^2\times 5^4}$}}
  {\ealpara{Połów każdą \ealkeytr{potęgę}:\\[0.3em]
   $2^4\times 3\times 5^2=16\times 3\times 25=1200$}{Halve each \ealkey{power}:\\[0.3em]
   $2^4\times 3\times 5^2=16\times 3\times 25=1200$}}

\qq[\Large][\large]{\ealpara{$n=2^3\times 3\times 5^2$\\[0.4em]
Znajdź najmniejszą liczbę całkowitą $k$, aby $nk$ było \ealkeytr{liczbą kwadratową}.}{$n=2^3\times 3\times 5^2$\\[0.4em]
Find the smallest whole number $k$ so that $nk$ is a \ealkey{square number}.}}
  {\ealpara{\ealkeytr{Potęgi} $2$ i $3$ są nieparzyste, więc każda potrzebuje jeszcze jednego:\\[0.3em]
   $k=2\times 3=6$, co daje $nk=2^4\times 3^2\times 5^2$.}{The \ealkey{powers} of $2$ and $3$ are odd, so each needs one more:\\[0.3em]
   $k=2\times 3=6$, giving $nk=2^4\times 3^2\times 5^2$.}}

\qq[\Large][\large]{\ealpara{$504=2^3\times 3^2\times 7$\\[0.4em]
Znajdź najmniejszą liczbę, przez którą można pomnożyć $504$, aby otrzymać \ealkeytr{liczbę kwadratową}.}{$504=2^3\times 3^2\times 7$\\[0.4em]
Find the smallest number $504$ can be multiplied by to give a \ealkey{square number}.}}
  {\ealpara{\ealkeytr{Potęgi} $2$ i $7$ są nieparzyste:\\[0.3em]
   pomnóż przez $2\times 7=14$.\\[0.3em]
   $504\times 14=7056=84^2$}{The \ealkey{powers} of $2$ and $7$ are odd:\\[0.3em]
   multiply by $2\times 7=14$.\\[0.3em]
   $504\times 14=7056=84^2$}}

\qq[\Large][\large]{\ealpara{$1080=2^3\times 3^3\times 5$\\[0.4em]
Znajdź najmniejszą liczbę, przez którą można podzielić $1080$, aby otrzymać \ealkeytr{liczbę kwadratową}.}{$1080=2^3\times 3^3\times 5$\\[0.4em]
Find the smallest number $1080$ can be divided by to give a \ealkey{square number}.}}
  {\ealpara{Usuń jedną $2$, jedną $3$ i $5$:\\[0.3em]
   podziel przez $2\times 3\times 5=30$.\\[0.3em]
   $1080\div 30=36=6^2$}{Remove one $2$, one $3$ and the $5$:\\[0.3em]
   divide by $2\times 3\times 5=30$.\\[0.3em]
   $1080\div 30=36=6^2$}}

%================================================================
\section{Problem solving with HCF and LCM}
\begin{frame}
  \vfill
  \begin{center}
    \ealpara{Rozwiązywanie problemów z \ealkeytr{HCF} i \ealkeytr{LCM}}{Problem solving with \ealkey{HCF} and \ealkey{LCM}}
    \par\vspace{1em}
    {\usebeamercolor[fg]{structure}\Huge\bfseries Problem solving with \ealkey{HCF} and \ealkey{LCM}}
  \end{center}
  \vfill
\end{frame}
%================================================================

\tf[\Large]{\ealpara{Aby znaleźć, kiedy dwa powtarzające się zdarzenia nastąpią razem,
używasz \ealkeytr{LCM}}{To find when two repeating events next happen together,
you use the \ealkey{LCM}}}
  {\ealpara{\textbf{Prawda.} Szukasz pierwszego czasu, który jest \ealkeytr{wielokrotnością}
   obu odstępów.}{\textbf{True.} You are looking for the first time that is a \ealkey{multiple}
   of both gaps.}}

\tf[\Large]{\ealpara{Aby podzielić dwie kolekcje na największe identyczne paczki,
używasz \ealkeytr{LCM}}{To split two collections into the largest identical bags,
you use the \ealkey{LCM}}}
  {\ealpara{\textbf{Fałsz.} Użyj \ealkeytr{HCF} --- rozmiar paczki musi być \ealkeytr{dzielnikiem}
   obu ilości.}{\textbf{False.} Use the \ealkey{HCF} --- the bag size has to be a \ealkey{factor}
   of both amounts.}}

\tf[\Large]{\ealpara{Cięcie dwóch lin na równe kawałki, tak długie, jak to możliwe, bez
resztek, używa \ealkeytr{HCF}}{Cutting two ropes into equal pieces, as long as possible with
none left over, uses the \ealkey{HCF}}}
  {\ealpara{\textbf{Prawda.} Długość kawałka musi być \ealkeytr{dzielnikiem} obu długości lin.}{\textbf{True.} The piece length must be a \ealkey{factor} of both rope lengths.}}

\tf[\Large]{\ealpara{\ealkeytr{LCM} dwóch liczb nigdy nie jest mniejsza od żadnej z nich}{\ealkey{The LCM} of two numbers is never smaller than either of them}}
  {\ealpara{\textbf{Prawda.} Jest \ealkeytr{wielokrotnością} każdej liczby,
   więc jest co najmniej tak duża jak obie.}{\textbf{True.} It is a \ealkey{multiple} of each number,
   so it is at least as big as both.}}

\tf[\Large]{\ealpara{Dwa zdarzenia powtarzają się co $4$ minuty i co $6$ minut,\\[0.3em]
więc spotykają się co $24$ minuty}{Two events repeat every $4$ minutes and every $6$ minutes,\\[0.3em]
so they coincide every $24$ minutes}}
  {\ealpara{\textbf{Fałsz.} Spotykają się co $\mathrm{LCM}(4,6)=12$ minut.}{\textbf{False.} They coincide every $\mathrm{LCM}(4,6)=12$ minutes.}}

\qq[\Large][\Large]{\ealpara{Dwa dzwonki dzwonią razem teraz.\\[0.4em]
Jeden dzwoni co $4$ minuty, drugi co $6$ minut.\\[0.4em]
Kiedy następnym razem zadzwonią razem?}{Two bells ring together now.\\[0.4em]
One rings every $4$ minutes, the other every $6$ minutes.\\[0.4em]
When do they next ring together?}}
  {\ealpara{Wspólna \ealkeytr{wielokrotność}, więc użyj \ealkeytr{LCM}:\\[0.3em]
   $\mathrm{LCM}(4,6)=12$ --- za $12$ minut.}{Common \ealkey{multiple}, so use the \ealkey{LCM}:\\[0.3em]
   $\mathrm{LCM}(4,6)=12$ --- in $12$ minutes.}}

\qq[\Large][\Large]{\ealpara{Dwie wstążki, o długościach $18$\,cm i $24$\,cm, są cięte na
równe kawałki, tak długie, jak to możliwe, bez resztek.\\[0.4em]
Jak długa jest każda część?}{Two ribbons, $18$\,cm and $24$\,cm long, are cut into
equal pieces, as long as possible with none left over.\\[0.4em]
How long is each piece?}}
  {\ealpara{Długość musi być \ealkeytr{dzielnikiem} obu, więc użyj \ealkeytr{HCF}:\\[0.3em]
   $\mathrm{HCF}(18,24)=6$ --- każda część ma $6$\,cm.}{The length must be a \ealkey{factor} of both, so use the \ealkey{HCF}:\\[0.3em]
   $\mathrm{HCF}(18,24)=6$ --- each piece is $6$\,cm.}}

\qq[\Large][\large]{\ealpara{$36$ pisaków i $48$ ołówków jest rozdzielonych do identycznych paczek
bez resztek.\\[0.4em]
Jaka jest największa liczba paczek?}{$36$ pens and $48$ pencils are shared into identical packs
with nothing left over.\\[0.4em]
What is the greatest number of packs?}}
  {\ealpara{$36=2^2\times 3^2$ \qquad $48=2^4\times 3$\\[0.4em]
   $\mathrm{HCF}=2^2\times 3=12$ paczek,\\[0.3em]
   każda z $3$ pisakami i $4$ ołówkami.}{$36=2^2\times 3^2$ \qquad $48=2^4\times 3$\\[0.4em]
   $\mathrm{HCF}=2^2\times 3=12$ packs,\\[0.3em]
   each with $3$ pens and $4$ pencils.}}

\qq[\Large][\Large]{\ealpara{Dwie latarnie morskie błyskają razem o $8$pm.\\[0.4em]
Jedna błyska co $15$ sekund, druga co $18$ sekund.\\[0.4em]
Kiedy następnym razem błysną razem?}{Two lighthouses flash together at $8$pm.\\[0.4em]
One flashes every $15$ seconds, the other every $18$ seconds.\\[0.4em]
When do they next flash together?}}
  {\ealpara{$\mathrm{LCM}(15,18)=90$ sekund\\[0.3em]
   O $8${:}$01${:}$30$pm.}{$\mathrm{LCM}(15,18)=90$ seconds\\[0.3em]
   At $8${:}$01${:}$30$pm.}}

\qq[\Large][\Large]{\ealpara{Prostokątne płytki mają wymiary $12$\,cm na $18$\,cm.\\[0.4em]
Jaki jest najmniejszy kwadrat, który można nimi dokładnie wyłożyć?}{Rectangular tiles measure $12$\,cm by $18$\,cm.\\[0.4em]
What is the smallest square they can tile exactly?}}
  {\ealpara{Bok musi być \ealkeytr{wielokrotnością} obu, więc użyj \ealkeytr{LCM}:\\[0.3em]
   $\mathrm{LCM}(12,18)=36$\\[0.3em]
   Kwadrat $36$\,cm na $36$\,cm.}{The side must be a \ealkey{multiple} of both, so use the \ealkey{LCM}:\\[0.3em]
   $\mathrm{LCM}(12,18)=36$\\[0.3em]
   A $36$\,cm by $36$\,cm square.}}

\qq[\Large][\Large]{\ealpara{Autobus odjeżdża co $24$ minuty, a tramwaj co
$36$ minut.\\[0.4em]
Oba odjeżdżają o $9${:}$00$. Kiedy następnym razem odjadą razem?}{A bus leaves every $24$ minutes and a tram every
$36$ minutes.\\[0.4em]
Both leave at $9${:}$00$. When do they next leave together?}}
  {\ealpara{$\mathrm{LCM}(24,36)=72$ minuty\\[0.3em]
   O $10${:}$12$.}{$\mathrm{LCM}(24,36)=72$ minutes\\[0.3em]
   At $10${:}$12$.}}

\qq[\Large][\large]{\ealpara{Dwa koła zębate, mające $18$ i $24$ zębów, zaczynają zazębione na
zaznaczonym zębie.\\[0.4em]
Ile obrotów wykona każde, zanim znaczniki znowu się spotkają?}{Two cogs, with $18$ and $24$ teeth, start meshed at a
marked tooth.\\[0.4em]
How many turns does each make before the marks meet again?}}
  {\ealpara{$\mathrm{LCM}(18,24)=72$ zębów musi przejść.\\[0.4em]
   $72\div 18=4$ obroty małego koła,\\[0.3em]
   $72\div 24=3$ obroty dużego koła.}{$\mathrm{LCM}(18,24)=72$ teeth must pass.\\[0.4em]
   $72\div 18=4$ turns of the small cog,\\[0.3em]
   $72\div 24=3$ turns of the large cog.}}

\qq[\Large][\large]{\ealpara{Podłoga ma wymiary $360$\,cm na $504$\,cm.\\[0.4em]
Jest wyłożona identycznymi kwadratowymi płytkami, tak dużymi, jak to możliwe.\\[0.4em]
Znajdź rozmiar płytki i liczbę płytek.}{A floor measures $360$\,cm by $504$\,cm.\\[0.4em]
It is tiled with identical square tiles, as large as possible.\\[0.4em]
Find the tile size and the number of tiles.}}
  {\ealpara{$360=2^3\times 3^2\times 5$ \qquad $504=2^3\times 3^2\times 7$\\[0.4em]
   $\mathrm{HCF}=2^3\times 3^2=72$, więc płytki mają $72$\,cm kwadratowe.\\[0.3em]
   $5\times 7=35$ płytek.}{$360=2^3\times 3^2\times 5$ \qquad $504=2^3\times 3^2\times 7$\\[0.4em]
   $\mathrm{HCF}=2^3\times 3^2=72$, so the tiles are $72$\,cm square.\\[0.3em]
   $5\times 7=35$ tiles.}}

\qq[\Large][\large]{\ealpara{Trzech biegaczy potrzebuje $60$\,s, $72$\,s i $90$\,s na okrążenie.\\[0.4em]
Zaczynają razem. Kiedy następnym razem wszyscy będą na linii startu?}{Three runners take $60$\,s, $72$\,s and $90$\,s per lap.\\[0.4em]
They start together. When are they next all at the start line?}}
  {\ealpara{$60=2^2\times 3\times 5$ \quad $72=2^3\times 3^2$ \quad
   $90=2\times 3^2\times 5$\\[0.4em]
   $\mathrm{LCM}=2^3\times 3^2\times 5=360$\,s\\[0.3em]
   Po $6$ minutach.}{$60=2^2\times 3\times 5$ \quad $72=2^3\times 3^2$ \quad
   $90=2\times 3^2\times 5$\\[0.4em]
   $\mathrm{LCM}=2^3\times 3^2\times 5=360$\,s\\[0.3em]
   After $6$ minutes.}}

\qq[\Large][\large]{\ealpara{Hot dogi są w paczkach po $10$, bułki w paczkach po $8$, a
saszetki sosu w paczkach po $12$.\\[0.4em]
Jaka jest najmniejsza liczba pełnych hot dogów, które Sam może zrobić, nie marnując
niczego, i ile paczek każdego produktu to wymaga?}{Hot dogs come in packs of $10$, buns in packs of $8$ and
sauce sachets in packs of $12$.\\[0.4em]
What is the smallest number of complete hot dogs Sam can make with nothing
left over, and how many packs of each is that?}}
  {\ealpara{$\mathrm{LCM}(10,8,12)=2^3\times 3\times 5=120$\\[0.4em]
   $120$ hot dogów: $12$ paczek hot dogów,\\[0.3em]
   $15$ paczek bułek i $10$ paczek saszetek.}{$\mathrm{LCM}(10,8,12)=2^3\times 3\times 5=120$\\[0.4em]
   $120$ hot dogs: $12$ packs of hot dogs,\\[0.3em]
   $15$ packs of buns and $10$ packs of sachets.}}

\end{document}